%% -*- coding:utf-8 -*-
\chapter{Passiv}
\label{chap-case}

%\if0

Dieses Kapitel behandelt das Passiv. Das Passiv wird üblicherweise als Unterdrückung des
Subjekts analysiert. Bevor ich jedoch eine Analyse entwickeln kann, muss ich fragen, was ein
Subjekt überhaupt ausmacht. Das ist eine Frage, der ganze Sammelbände und Dissertationen gewidmet
sind, und so bescheiden ich bin, werde ich versuchen, zumindest für die \ili{germanisch}en Sprachen
eine Antwort zu geben. Wie wir sehen werden, ist die Situation in Sprachen wie dem \ili{Dänisch}en,
\ili{Englisch}en und \ili{Deutsch}en recht klar, aber über das \ili{Isländisch}e gibt es spannende
Fakten zu entdecken.

\section{Das Phänomen}

\subsection{Subjekte und andere Subjekte}
\label{sec-subj-properties}
\label{sec-icelandic-quirky-subj}

Die\is{Subjekt|(} Situation in Sprachen wie dem \ili{Dänisch}en, \ili{Englisch}en und \ili{Deutsch}en ist recht klar. Viele
Autoren nehmen beispielsweise an, dass nicht-prädikative NPs im Nominativ Subjekte im \ili{Deutsch}en sind. So ist \emph{der
  Delphin} das Subjekt der Sätze in (\mex{1}):
\eal
\ex 
Der Delphin lacht.
\ex 
Der Delphin hilft dem Kind.
\ex 
Der Delphin gibt  ihr einen Ball.
\zl
Die Beschränkung auf nicht-prädikative NPs ist nötig, da wir andernfalls annehmen müssten, dass beide
NPs in (\mex{1}) Subjekte sind; aber \emph{ein Lügner} ist eine prädikative Phrase und nur
\emph{der Mann} ist das Subjekt.
\ea
Der Mann ist ein Lügner.
\z
Außerdem werden bestimmte satzwertige Argumente als Subjekte behandelt.

Genitive und Dative wie in (\mex{1}) werden im \ili{Deutsch}en nicht zu den Subjekten gezählt.
\eal
\ex 
Ihrer       wurde gedacht.
\ex 
Ihm       wurde  geholfen.
\zl

\noindent
Interessanterweise wurde die Frage, ob Genitive und Dative wie die in (\mex{0}) Subjekte sind, für
die SVO-Sprache \ili{Isländisch} von Forschern, die der Arbeit von
\citet*{ZMT85a} folgen, ganz anders beantwortet. Obwohl die Sätze in (\mex{1}) wie die in (\mex{-1}) aussehen, wird für das genitivische und
das dativische Element in (\mex{1}a) und (\mex{1}b) behauptet, dass sie Subjekte seien.
\eal
\label{ex-subject-icelandic-passive-v2}
\ex 
\gll Hennar var saknað.\\
     sie.\SG.\GEN{} war vermisst\\\icelandic
\ex 
\gll Þeim            var hjálpað.\\
     sie.\PL.\DAT{} war geholfen\\
\zl
Da das \ili{Isländisch}e eine V2-Sprache ist, hilft uns die Konstituentenabfolge in solchen einfachen Sätzen nicht dabei zu
bestimmen, ob \emph{hennar} und \emph{Þeim} Subjekte sind oder nicht. Diese Elemente sind
vorangestellt, und da sowohl Subjekte als auch Objekte vorangestellt werden können, helfen uns die Sätze in (\mex{0}) nicht
dabei, die grammatische Funktion dieser Argumente zu bestimmen. \citet*{ZMT85a} haben jedoch argumentiert, dass
diese Elemente als Subjekte analysiert werden sollten, und haben eine Testbatterie bereitgestellt. Zu den Tests gehören aufwändigere
Stellungstests und die Weglassbarkeit in sogenannten Kontrollkonstruktionen. Ich wende mich nun diesen
Tests zu.

\subsubsection{Die Position von Subjekten in V2- und V1-Sätzen}

Der\is{Abfolge!V2|(}\is{Abfolge!V1|(} erste vorgeschlagene Test nutzt die Position von Konstituenten in V2-Sätzen, in denen ein
Nicht-Subjekt vorangestellt ist \citep*[Abschnitt~2.3]{ZMT85a}. Betrachten wir beispielsweise die folgenden Beispiele:

\eal
%% \ex[]{
%% \gll Refinn           skaut  Ólafur      með þessari byssu.\\
%%      den.Fuchs.\ACC{} schoss Olaf.\NOM{} mit diesem  Gewehr\\
%% }
% selbst erfunden, check
\ex[]{
\gll Með  þessari byssu   skaut Ólafur      refinn.\\
     mit diesem Gewehr schoss Olaf.\NOM{} der.Fuchs.\ACC{}\\\icelandic
}
\ex[*]{
\gll Með þessari byssu  skaut refinn         Ólafur.\\
     mit diesem Gewehr schoss der.Fuchs.\ACC{} Olaf.\NOM{}\\
}
\zl
Der Nominativ kann unmittelbar nach dem finiten Verb \emph{skaut} `schoss' wie in (\mex{0}a) auftreten, aber er
kann nicht rechts vom Akkusativ wie in (\mex{0}b) stehen.

Der zweite Test nutzt \emph{w}-Fragen und prüft die Position des Subjekts in Bezug auf das
Objekt und auf nicht-finite Verben:
\eal
\ex[]{
\gll Hvenær hafði Sigga        hjálpað Haraldi?\\
     wann hat Sigga.\NOM{} geholfen Harald.\DAT{}\\\icelandic
}
\ex[*]{
\gll Hvenær hafði Haraldi       Sigga        hjálpað?\\
     wann hat Harald.\DAT{} Sigga.\NOM{} geholfen\\
}
\zl
Das Objekt muss dem Partizip \emph{hjálpað} folgen wie in (\mex{0}a), und das Subjekt folgt
unmittelbar dem finiten Verb. Beispiele mit dem Objekt vor dem Subjekt wie in (\mex{0}b) sind
ungrammatisch. Das Dativobjekt kann vorangestellt werden, muss dann aber in Erststellung
links vom finiten Verb realisiert werden, nicht rechts davon:
\ea[]{
\gll Haraldi       hafði Sigga        aldrei hjálpað.\\
     Harald.\DAT{} hat Sigga.\NOM{} nie geholfen\\\icelandic
}
\z

\noindent
Dieselbe Situation findet sich bei Entscheidungsfragen:
\eal
\ex[]{
\gll Hafði Sigga        aldrei hjálpað Haraldi?\\
     hat Sigga.\NOM{} nie geholfen Harald.\DAT\\\icelandic
}
\ex[*]{
\gll Hafði Haraldi       Sigga        aldrei hjálpað?\\
     hat Harald.\DAT{} Sigga.\NOM{} nie geholfen\\
}
\zl

\noindent
\citet*[Abschnitt~2.3]{ZMT85a} hat beobachtet, dass bestimmte Dative ebenfalls in dieser postverbalen Position auftreten können:

\eal
\ex[]{
\gll Hefur henni      alltaf þótt    Ólafur      leiðinlegur?\\
     hat sie.\DAT{} immer gedacht Olaf.\NOM{} langweilig.\NOM{}\\\icelandic
\glt `Hat sie Olaf immer langweilig gefunden?'
}
\ex[]{
\gll Ólafur      hefur henni alltaf þótt leiðinlegur.\\
     Olaf.\NOM{} hat sie.\DAT{} immer gedacht langweilig.\NOM{}\\
\glt `Sie hat Olaf immer langweilig gefunden.'
}
\ex[*]{
\gll Hefur Ólafur      henni      alltaf þótt    leiðinlegur?\\
     hat Olaf.\NOM{} sie.\DAT{} immer gedacht langweilig.\NOM\\
}
\zl

\noindent
Das \ili{deutsch}e Äquivalent wäre der Satz in (\mex{1}):
{\judgewidth{??}
\ea[??]{
Mich     dünkt  der Film langweilig.
}
\z}

\noindent
\emph{dünkt} ist allerdings archaisch und wird üblicherweise mit einem \emph{dass}-Satz verwendet -- wenn es überhaupt
verwendet wird. Es gibt aber ein nicht-archaisches Verb mit ähnlicher Form:
\ea
Mir scheint der Mann langweilig.
\z
Der Experiencer von \emph{scheinen} wird mit dem Dativ ausgedrückt, während das Subjekt des
eingebetteten Prädikats \emph{langweilig} im Nominativ steht.
\is{Abfolge!V2|)}\is{Abfolge!V1|)}

\subsubsection{Subjekte in Kontrollkonstruktionen}
\label{sec-subject-control}

\citet*[Abschnitt~2.7]{ZMT85a}\is{Kontrolle|(} diskutieren Kontrollstrukturen, in denen das Subjekt des eingebetteten Verbs
nicht ausgedrückt ist. (\mex{1}a) zeigt ein Beispiel für normale Kontrolle, in dem das Subjekt des Matrixverbs
\emph{vonast} `hoffen' sich auf denselben Diskursreferenten bezieht wie das Subjekt des eingebetteten Verbs \emph{fara} `gehen'. (\mex{1}b) ist ein Beispiel für sogenannte arbiträre Kontrolle. Bei arbiträrer Kontrolle gibt es
kein vom infinitregierenden Kopf abhängiges Element, das sich auf denselben Diskursreferenten
bezieht wie das Subjekt des Infinitivs. Das nicht ausgedrückte Subjekt entspricht einem Pronomen \emph{man}, das
generisch verwendet wird. In Beispiel (\mex{1}b) selegiert \emph{óvenjulegt} `ungewöhnlich' kein Argument,
das sich auf denselben Referenten bezieht wie das Subjekt von \emph{fara} `gehen'. Das Subjekt von \emph{að fara
  heim snemma} `früh nach Hause gehen' ist nicht ausgedrückt, wird aber als das Indefinitpronomen \emph{man} verstanden.
\eal
\ex
\gll Ég  vonast til að fara heim.\\
     ich hoffe um zu gehen heim\\\icelandic
\glt `Ich hoffe, nach Hause zu gehen.'
\ex
\gll Að fara heim snemma er óvenjulegt.\\
     zu gehen heim früh ist ungewöhnlich\\
\glt `Es ist ungewöhnlich, früh nach Hause zu gehen.'
\zl

Nun kann man beobachten, dass das \ili{Isländisch}e Verben erlaubt, die in solchen Kontrollkonstruktionen keinen Nominativ nehmen. Ein Beispiel ist \emph{vantar} (`fehlen'), das zwei Akkusative statt eines
Nominativs und eines Akkusativs nimmt:
\ea
\gll Mig      vantar peninga.\\
     ich.\ACC{} fehlt Geld.\ACC\\\icelandic
\z
(\mex{1}) zeigt, dass dieses Verb unter \emph{vonast} (`hoffen') eingebettet werden kann:
\ea
\gll Ég vonast til að vanta ekki peninga.\\
     ich hoffe um zu fehlen nicht Geld.\ACC\\\icelandic
\glt `Ich hoffe, dass mir kein Geld fehlt.'
\z

Dies sollte mit dem \ili{Deutsch}en verglichen werden:
\eal
\ex[]{
Mir fehlt kein Geld.
}
\ex[*]{
Ich hoffe, kein         Geld  zu fehlen.
}
\zl

\noindent
Die Frage am Anfang dieses Abschnitts war, ob die Dative und Genitive in Sätzen
wie (\ref{ex-subject-icelandic-passive-v2}), die hier als (\mex{1}) wiederholt sind, Subjekte sind oder nicht.
\eal
\ex 
\gll Hennar     var saknað.\\
     sie.\GEN{} war vermisst\\\icelandic
\glt `Sie wurde vermisst.'
\ex 
\gll Þeim        var hjálpað.\\
     sie.\DAT{} war geholfen\\
\glt `Ihnen wurde geholfen.'
\zl

\noindent
Wir können nun die Tests verwenden, um diese Frage zu beantworten: Der Dativ ist in der Frage in (\mex{1}) unmittelbar rechts vom finiten
Verb und steht somit in Subjektposition.

\eal
\ex 
\gll Var hennar saknað?\\
     war sie.\GEN{} vermisst\\\icelandic
\glt `Wurde sie vermisst?'
\ex
\gll Var þeim       hjálpað?\\
     war sie.\DAT{} geholfen\\
\glt `Wurde ihnen geholfen?'
\zl

% \ea
% \gll Var honum     aldrei hjálpað af foreldrum sinum?\\
%      was he.\DAT{} never  helped   by parents   his\\
% \glt `Did his parents never help him?'
% \z

\noindent
Ebenso folgt der Dativ dem finiten Verb im V2-Satz in (\mex{1}):
\ea
\gll Í prófinu  var honum víst hjálpað.\\
     in der.Prüfung war er.\DAT{} offenbar geholfen\\\icelandic
\glt `Offenbar wurde ihm in der Prüfung geholfen.'
\z

\noindent
Außerdem können diese Dative in Kontrollkonstruktionen weggelassen werden, wie die Beispiele in (\mex{1}) zeigen:
\eal
\ex
\gll Ég vonast til að verða hjálpað.\\
     ich hoffe um zu werden geholfen\\\icelandic
\ex
\gll Að vera hjálpað i prófinu er óleyfilegt.\\
     zu sein geholfen in der.Prüfung ist un.erlaubt\\
\glt `Es ist nicht erlaubt, dass einem in der Prüfung geholfen wird.'
\zl
Dies sollte mit dem \ili{Deutsch}en verglichen werden: Während Verben wie \emph{unterstützen}, die einen
Nominativ und einen Akkusativ regieren, in solchen Kontrollkonstruktionen auftreten können, sind Verben wie \emph{helfen}, die einen Nominativ und einen Dativ nehmen, in dieser Konstruktion ausgeschlossen:
\eal
\ex 
dass jemand ihm hilft
\ex
dass jemand ihn unterstützt
\ex
dass ihm geholfen wird
\ex
dass er unterstützt wird
\zl

\eal
\ex[]{
Ich hoffe unterstützt zu werden.
}
\ex[*]{
Ich hoffe geholfen zu werden.
}
\zl
Das Dativobjekt kann in solchen Kontrollkonstruktionen nicht weggelassen werden, wie (\mex{0}b) zeigt. Die einzige Möglichkeit,
unter \emph{hoffen} in einem Infinitivsatz ein Passiv zu realisieren, ist das Dativpassiv mit
\emph{erhalten}/""\emph{bekommen}/""\emph{kriegen}. Das Dativpassiv kann ein Dativobjekt in ein
nominativisches Subjekt verwandeln:\footnote{
  Nicht alle Sprecher akzeptieren solche Sätze. Wir kommen weiter unten bei der Diskussion der Beispiele
  (\ref{ex-helfen-with-dative-passive}) und (\ref{ex-helfen-with-dative-passive-corpus}) darauf zurück.
}
\ea[]{
Aicke        bekommt geholfen.
}
\z
\largerpage[-1]
Da das Objekt von \emph{helfen} dann im Nominativ steht und somit zweifellos ein Subjekt im \ili{Deutsch}en ist, überrascht es nicht, dass es in Kontrollkonstruktionen wie (\mex{1}) weggelassen werden kann:
\ea[]{
Ich hoffe hier geholfen zu bekommen.\footnote{    
\url{http://www.photovoltaikforum.com/sds-allgemein-ueber-solar-log-f38/solarlog-1000-mit-wifi-anschliesen-t96371.html}. 10.01.2014
}

}
\z
\is{Subjekt|)}\is{Kontrolle|)}

\subsection{Vergleich zwischen Deutsch, Dänisch, Englisch und Isländisch}

In den folgenden Unterabschnitten vergleiche ich mehrere Dimensionen, in denen die \ili{germanisch}en Sprachen
variieren:
\begin{itemize}
\item Das \ili{Dänisch}e und das \ili{Isländisch}e haben ein morphologisches Passiv; das \ili{Englisch}e und das \ili{Deutsch}e nicht.
% ZMT85a:443 \ili{Icelandic} also has some morphologically'middle' forms in the suffix -st, someof which
% have a passive meaning, as illustrated in (3):5

\item Das \ili{Deutsch}e und das \ili{Isländisch}e erlauben subjektlose Konstruktionen; das \ili{Dänisch}e und das \ili{Englisch}e nicht.

\item Das \ili{Dänisch}e, das \ili{Deutsch}e und das \ili{Isländisch}e erlauben unpersönliche Passive; das \ili{Englisch}e nicht.

\item Das \ili{Dänisch}e und das \ili{Isländisch}e erlauben, dass beide Objekte zum Subjekt befördert werden; das \ili{Englisch}e und das \ili{Deutsch}e nicht.

\item Das \ili{Deutsch}e hat das Fernpassiv, das \ili{Dänisch}e das komplexe Passiv, und das \ili{Dänisch}e und das \ili{Englisch}e
  haben das reportive Passiv.
\end{itemize}

\subsubsection{Morphologische und analytische Formen}

Das \ili{Dänisch}e hat ein morphologisches\is{Morphologie} Passiv\is{Passiv!morphologisches}. Es wird durch Anfügen des Suffixes \suffix{s} an das Verb gebildet, und es
gibt Formen für das Präsens (\ref{ex-laeses}) und das Präteritum (\ref{ex-laestes}):
\eal
\ex[]{\label{ex-laeseract}
\gll Peter læser avisen.\\
     Peter liest Zeitung.\textsc{def}\\\danish
\glt `Peter liest die Zeitung.'}
\ex[]{\label{ex-laeses}
\gll Avisen              læses af Peter.\\
     Zeitung.\textsc{def} lesen.\textsc{pres}.\textsc{pass} von Peter\\
\glt `Die Zeitung wird von Peter gelesen.'}
\ex[]{\label{ex-laestes}
\gll Avisen              læstes af Peter.\\
     Zeitung.\textsc{def} lesen.\textsc{past}.\textsc{pass} von Peter\\
\glt `Die Zeitung wurde von Peter gelesen.'}
\zl
Wie die Beispiele in (\mex{1}) zeigen, ist die \emph{af}-Phrase nicht notwendig:
\eal
\ex 
\gll Avisen        læses           hver   dag.\\
     Zeitung.\textsc{def} liest jeden Tag\\\danish
\glt `Die Zeitung wird jeden Tag gelesen.'
\ex 
\gll Avisen        læstes           hver   dag.\\
     Zeitung.\textsc{def} las jeden Tag\\\danish
\glt `Die Zeitung wurde jeden Tag gelesen.'
\zl

\noindent
Das \ili{Dänisch}e hat auch eine analytische Form mit \emph{blive} `werden' plus Partizip~II:
\ea
\gll Avisen                 bliver læst af Peter.\\
     Zeitung.\textsc{def} wird gelesen von Peter\\\danish
\glt `Die Zeitung wird von Peter gelesen.'
\z
Das morphologische Passiv kann auch auf Infinitive angewendet werden:
\ea
\gll Avisen skal læses hver dag.\\
     Zeitung.def muss lesen.\textsc{inf}.\textsc{pass} jeden Tag\\\danish
\glt `Die Zeitung muss jeden Tag gelesen werden.'
\z

\noindent
Das \ili{Englisch}e und das \ili{Deutsch}e haben nur die analytische Variante, also die Variante mit einem
Hilfsverb:
\eal
\ex The paper was read.
\ex 
Der        Aufsatz wurde  gelesen.
\zl    

\subsubsection{Persönliches und unpersönliches Passiv}
\label{sec-impersonal-passive-phen}

Alle\is{Passiv!persönliches|(}\is{Passiv!unpersönliches|(} betrachteten Sprachen erlauben die Beförderung eines Akkusativobjekts zum Subjekt, ein
Beispiel dafür ist in (\mex{1}) gegeben.
\eal
\ex
\longexampleandlanguage{
Angehörige haben den Verdächtigen zuletzt am Montag gesehen.
\ex 
Der        Verdächtige wurde  zuletzt am Montag gesehen.
\zl

\noindent
Wie die folgenden Beispiele zeigen, kann das Subjekt ein S oder eine VP sein:\is{Subjektsatz}
% \todostefan{show that these are really
%   subjects. We have V2 here, could be objects.}

\eal
\ex
\gll At regeringen træder tilbage, bliver påstået.\\
     dass Regierung.\textsc{def} tritt \textsc{part} wird behauptet\\\danish
\glt `Es wird behauptet, dass die Regierung zurücktritt.'
\ex
\gll At reparere bilen, bliver forsøgt.\\
     zu reparieren Auto.\textsc{def} wird versucht\\
\glt `Es wird versucht, das Auto zu reparieren.'
\zl

\noindent
Zusätzlich zu solchen persönlichen Passiven haben das \ili{Dänisch}e, das \ili{Deutsch}e und das \ili{Isländisch}e unpersönliche
Passive.\footnote{
Die Bezeichnungen ``persönliches'' und ``unpersönliches Passiv'' sind irreführend, da beide Passive die Eigenschaft teilen, das Subjekt zu degradieren. Sogenannte persönliche Passive können belebte oder unbelebte Subjekte haben:
\ea 
Der        Diamant wurde  zuletzt am Montag gesehen.
\z
Die große \ili{deutsch}e Grammatik des Instituts für Deutsche Sprache hat versucht, die neuen Begriffe \emph{Zweitakt-Passiv} und
\emph{Eintakt-Passiv} zu etablieren \citep[\page 1793]{Zifonun97b}. Die erste Phase ist dabei die Unterdrückung des Subjekts und
die zweite Phase die Beförderung des Akkusativobjekts zum Subjekt bei denjenigen Verben, die einen
Akkusativ regieren. Diese Begriffe sind im Prinzip zwar angemessener, ich verwende sie hier aber nicht, da die
im Folgenden vorgeschlagene Analyse beide Passive einheitlich behandelt: Sie unterdrückt lediglich das
Subjekt. Persönliche und unpersönliche Passive werden gleich analysiert. Der Unterschied ergibt sich aus
Unterschieden in der Kasuszuweisung. In Ermangelung besserer Begriffe verwende ich weiterhin die Begriffe persönliches und
unpersönliches Passiv.
} Da das \ili{Deutsch}e kein Subjekt verlangt, sind unpersönliche Passive wie (\mex{1}) zu erwarten:

\ea
weil    noch  getanzt wurde
\z

\noindent
Die folgenden beiden Beispiele aus dem \ili{Isländisch}en zeigen, dass das \ili{Isländisch}e ebenfalls unpersönliche Konstruktionen hat \citep[\page 264]{Thrainsson2007a-u}:
\eal
\ex 
\gll Oft var   talað      um   þennan mann.\\
     oft war geredet über diesen Mann.\ACC.\SG.\M\\\icelandic
\glt `Über diesen Mann wurde oft geredet.'
\ex
\gll Aldrei hefur verið    sofið      í  þessu  rúmi.\\
     nie hat gewesen geschlafen in diesem Bett.\DAT\\
\glt `In diesem Bett wurde noch nie geschlafen.'
\zl

\noindent
Das \ili{Dänisch}e erlaubt ebenfalls unpersönliche Passive, unterscheidet sich aber von den bisher besprochenen Sprachen darin, dass es ein expletives\is{Pronomen!expletives} Subjekt verlangt:
\eal
\ex\label{ex-at-der-bliver-danset} 
\gll at der bliver danset\\
     dass \textsc{expl} \AUX{} getanzt\\\danish
\glt `dass getanzt wird'
\ex
\gll at der danses\\
     dass \textsc{expl} tanzen.\textsc{pres}.\textsc{pass}\\
\glt `dass getanzt wird'
\ex[*]{ 
\gll Bliver danset?\\
     \AUX{} getanzt\\
}
\ex[*]{
\gll Danses?\\
     tanzen.\textsc{pass}\\
}
\zl
Das \ili{Dänisch}e ist also wie das \ili{Englisch}e darin, dass es immer ein Subjekt verlangt; während diese Beschränkung im \ili{Englisch}en aber zur
Unmöglichkeit unpersönlicher Passive führt, hat das \ili{Dänisch}e eine Lösung für das Subjektproblem gefunden, indem es
ein Expletivum einfügt. Interessanterweise ist diese Lösung Teil einer allgemeineren ``Strategie'' der
Expletiveinfügung. Das Dänische erlaubt die Einfügung von Expletiva bei verschiedenen Verbtypen, sofern
sie kein direktes Objekt haben. So können Passivverben und intransitive Verben mit
Expletiva kombiniert werden.
\eal
\ex 
\gll at    der    ikke går   en mand på gaden\\
     dass \expl{} nicht geht ein Mann auf Straße.\textsc{def}\\\danish
\ex 
\gll at   der     ikke kommer to  nye medarbejdere\\
     dass \expl{} nicht kommen zwei neue Mitarbeiter\\
\ex
\gll at   der     ikke venter nogle hårde forhandlinger\\
     dass \expl{} nicht warten einige harte Verhandlungen\\
\zl
Die Beispiele in (\mex{0}) sind eingebettete Sätze mit Negationen, was zeigt, dass das Expletivum nicht
in einem V2-Satz vorangestellt ist, sondern dass wir hier tatsächlich ein Expletivum in Subjektposition haben.

Expletiva\is{Pronomen!expletives} sind in unpersönlichen Konstruktionen des \ili{Deutsch}en ausgeschlossen:
\nocite{MOe2011a}
\ea[*]{
weil    es noch  gearbeitet wurde
}
\z
\is{Passiv!persönliches|)}\is{Passiv!unpersönliches|)}

%% The examples in (\ref{ex-gearbeitet-wurde}) and (\ref{ex-bliver-arbejder}) show passives of
%% mono-valent verbs but of course bi-valent intransitive verbs like the \ili{German} \emph{denken} (`think')
%% and \ili{Danish} \emph{passe} (`take care of') also form impersonal passives:
%% \ea
%% \gll dass an die Männer gedacht wurde\\
%%      that \textsc{prep} the men thought was\\
%% \glt `that one thought about the men'
%% \z
%% \eal
%% \label{ex-impersonal-passive-pp}
%% \ex
%% \gll Der passes på børnene.\\
%%      \textsc{expl} take.care.of.\textsc{pres}.\textsc{pass} on children.\textsc{def}\\
%% \glt `Somebody takes care of the children.'
%% \ex
%% \gll Der bliver passet  på børnene.\\
%%      \textsc{expl} is taken.care.of on children.\textsc{def}\\
%% \glt `Somebody takes care of the children.'
%% \zl

\subsubsection{Beförderung des primären und sekundären Objekts}

\largerpage
Das \ili{Englisch}e und das \ili{Deutsch}e erlauben nur die Beförderung eines der Objekte eines ditransitiven Verbs. (\mex{1})
zeigt, dass das Akkusativobjekt als Subjekt realisiert werden kann, das Dativobjekt jedoch nicht:
\eal
\ex[]{
\longexampleandlanguage{
weil    der        Mann dem Jungen den Ball schenkt
}
\ex[]{
weil    dem        Jungen der        Ball geschenkt wurde
}
\ex[*]{
weil der Junge den Ball geschenkt wurde
}
\zl

\noindent
Ebenso kann das \ili{Englisch}e das erste Objekt als Subjekt realisieren, das zweite Objekt kann aber nicht zum Subjekt befördert werden:
\eal
\ex[]{
because the man gave the child the ball
}
\ex[]{
because the child was given the ball
}
\ex[*]{
because the ball was given the child
}
\zl
Der informationsstrukturelle Effekt lässt sich allerdings mit einer anderen lexikalischen Variante von \emph{give}
erreichen. \emph{give} kann mit einem NP-Objekt und einer \emph{to}-PP statt mit zwei NPs verwendet werden, wie in
(\mex{1}a). Das erste Objekt des ditransitiven \emph{give} wird in (\mex{1}a) als PP realisiert, und das
zweite Objekt \emph{the ball} ist in (\mex{1}a) das erste Objekt. Diese Alternation wird auch Dativ-Shift\is{Dativ-Shift} genannt.
\eal
\ex because the man gave the ball to the child
\ex because the ball was given to the child
\zl
(\mex{0}b) ist die Passivvariante von (\mex{0}a). Wie in (\mex{-1}b) wird das primäre Objekt zum
Subjekt befördert.

Das \ili{Dänisch}e und das \ili{Isländisch}e unterscheiden sich vom \ili{Englisch}en und \ili{Deutsch}en. In den ersteren Sprachen können beide Objekte zum Subjekt befördert werden, ohne dass zuvor eine Alternation von Valenzrahmen wie der Dativ-Shift stattfindet.
\eal
\ex\label{ex-fordi-manden-giver-barnet-bolden} 
\gll fordi manden giver barnet bolden\\ 
     weil Mann.\textsc{def} gibt Kind.\textsc{def} Ball.\textsc{def}\\\danish
\glt `weil der Mann dem Kind den Ball gibt'
\ex\label{ex-child-was-given-ball-danish}
\gll fordi barnet bliver givet bolden\\ 
     weil Kind.\textsc{def} wird gegeben Ball.\textsc{def}\\
\glt `weil dem Kind der Ball gegeben wird'
\ex\label{ex-ball-was-given-child-danish}
\gll fordi bolden bliver givet barnet\\ 
     weil Ball.\textsc{def} wird gegeben Kind.\textsc{def}\\
\glt `weil der Ball dem Kind gegeben wird'
\zl

\noindent
Man könnte annehmen, dass immer das erste Objekt (das primäre Objekt) zum Subjekt befördert wird
und dass das \ili{Dänisch}e keine Abfolge der Objekte hat, sodass beide Objekte gleich prominent sind
und zum Subjekt befördert werden können. Das \ili{Moro} ist eine Sprache, der solche Eigenschaften zugeschrieben werden
\citep{AMM2017a-u}. Das \ili{Dänisch}e unterscheidet sich jedoch vom Moro darin, dass die Abfolge der Objekte in Sätzen
eindeutig festgelegt ist: Während (\mex{0}a) möglich ist, ist die umgekehrte Abfolge der Objekte ungrammatisch, wie
(\mex{1}) zeigt.
\ea[*]{
\gll fordi   manden           giver bolden            barnet \\
     weil Mann.\textsc{def} gibt Ball.\textsc{def} Kind.\textsc{def}\\\danish
}
\z

\largerpage
\noindent
Was das \ili{Isländisch}e betrifft, merken \citet*[\page 460]{ZMT85a} an, dass neben der Möglichkeit,
den Akkusativ zum nominativischen Subjekt zu befördern, der Dativ zu einem quirky Subjekt werden kann:
\ea
\label{ex-dat-subj-passive-ditransitive-icelandic}
\gll Konunginum voru gefnar ambáttir.\\
     der.König.\DAT{} waren gegeben.\F.\PL{} Mägde.\NOM.\F.\PL\\\icelandic
\glt `Dem König wurden Sklavinnen gegeben.'
\z
Die Struktur von (\mex{0}) ist in (\mex{1}) skizziert:
\ea
\label{sub-aux-v-o}
{}[S$_i$ Aux \_$_i$ V O] 
\z
Da der Nominativ nach dem Partizip angeordnet ist, kann er nicht das Subjekt sein, was impliziert, dass
das vorangestellte dativische Element das Subjekt ist.

Alternativ wird das Akkusativobjekt zum nominativischen Subjekt befördert:
\ea
\label{ex-nom-subj-passive-ditransitive-icelandic}
\gll Ambáttin var gefin konunginum.\\
     die.Magd.\NOM.\SG{} \AUX{} gegeben.\F.\SG{} der.König.\DAT\\\icelandic
\glt `Die Sklavin wurde dem König gegeben.'
\z
Auch dieser Satz hat die Struktur in (\ref{sub-aux-v-o}).

Um zu zeigen, dass der Dativ in
(\ref{ex-dat-subj-passive-ditransitive-icelandic}) tatsächlich zum Subjekt befördert wird und dass der Akkusativ in
(\ref{ex-nom-subj-passive-ditransitive-icelandic}) zum Subjekt befördert wird, wenden \citet*[\page 460]{ZMT85a} eine Batterie von
Tests an. Ich gebe hier nur die V2-Beispiele mit einem Adjunkt in Erststellung, die Fragen und die
Kontrollstrukturen an. Die Beispiele in (\mex{1}) und (\mex{2}) zeigen, dass die obigen Sätze
tatsächlich die Struktur in (\ref{sub-aux-v-o}) haben. Die erste Position in (\mex{1}) ist mit einem
Adjunkt gefüllt, was bedeutet, dass das Subjekt in Subjektposition verbleibt, und zeigt somit, dass der Dativ
\emph{konunginum} `der König' das Subjekt ist. Ebenso ist der Nominativ \emph{ambáttin} `die
Sklavin' das Subjekt in (\mex{1}b).

\eal
\ex
\gll Um veturinn voru konunginum gefnar ambáttir.\\
     in dem.Winter \AUX{} der.König.\DAT{} gegeben Sklavinnen.\NOM\\\icelandic
\glt `Im Winter wurden dem König Sklavinnen gegeben.'
\ex
\gll Um veturinn var ambáttin gefin konunginum.\\
     in dem.Winter \AUX{} die.Sklavin.\NOM{} gegeben der.König.\NOM\\
\glt `Im Winter wurde die Sklavin dem König gegeben.'
\zl
Die Fragen in (\mex{1}) sind ein weiterer Beleg. Die erste Position ist nicht gefüllt, und der Dativ
in (\mex{1}a) und der Nominativ in (\mex{1}b) werden unmittelbar nach dem finiten Verb realisiert.
\eal
\ex\label{ex-were-the-king-given-the-slaves}
\gll Voru konunginum gefnar ambáttir?\\
     \AUX{} der.König.\DAT{} gegeben Sklavinnen.\NOM{}\\\icelandic
\glt `Wurden dem König Sklavinnen gegeben?'
\ex\label{ex-were-the-slaves-given-the-king}
\gll Var ambáttin gefin konunginum?\\
     \AUX{} die.Sklavin.\NOM{} gegeben der.König.\DAT\\
\glt `Wurde die Sklavin dem König gegeben?'
\zl

\noindent
(\mex{1}) zeigt die entsprechenden Kontrollbeispiele:
\eal
\ex
\gll Að vera gefnar ambáttir var mikill heiður.\\
     zu \AUX{} gegeben Sklavinnen.\NOM{} war große Ehre\\\icelandic
\glt `Sklavinnen gegeben zu bekommen, war eine große Ehre.'
\ex
\gll Að vera gefin konunginum olli miklum vonbrigðum.\\
     zu \AUX{} gegeben der.König.\DAT{} verursachte große Enttäuschung\\
\glt `Dem König gegeben zu werden, verursachte große Enttäuschung.'

%% \ex
%% \gll Ambáttin vonast til að verða gefin konunginum.\\
%%      the.slave.\NOM{} hopes for to be given the.king.\DAT\\
%% \glt `The slave hopes to be given to the king.'
\zl
In (\mex{0}a) ist der Dativ nicht ausgedrückt und in (\mex{0}b) ist der Nominativ ausgelassen. Dies zeigt,
dass sowohl das primäre als auch das sekundäre Objekt im \ili{Isländisch}en zum Subjekt befördert werden können, obwohl
das primäre Objekt im Dativ steht und sich der Kasus der NP in Passivbeispielen nicht zum Nominativ ändert.

% \if0

% \subsubection{Subjekt-Verb-Kongruenz?}

% \begin{itemize}
% \item Verben kongruieren mit dem Nominativelement.\\
%       Wenn es keins gibt, ist das Verb dritte Person Singular (Neutrum).

% \item In (\mex{1}) keine Kongruenz:
% \eal
% \ex 
% \gll Þeim       var hjalpað.\\
%      sie.\emph{\PL}.\DAT{} wurde geholfen\\
% \ex 
% \gll Hennar var saknað.\\
%      sie.\emph{\SG}.\GEN{} wurde vermisst\\
% \zl

% \item Der Dativ und der Genitiv sind aber dennoch Subjekte,\\
%       wie wir gleich sehen werden.

% \end{itemize}

% \fi

\section{Die Analyse}

\subsection{Struktureller und lexikalischer Kasus und das Kasusprinzip}
\label{sec-struk-lex-kas}
\label{sec-struc-lex-kas}

Für\is{Kasus|(}\is{Kasus!struktureller|(}\is{Kasus!lexikalischer|(} die Analyse des Passivs ist es nützlich, zwischen strukturellem und lexikalischem
Kasus zu unterscheiden. Struktureller Kasus ist Kasus, der von der syntaktischen Struktur abhängt, in der Argumente
realisiert werden, während lexikalischer Kasus Kasus ist, der unabhängig von der syntaktischen Umgebung konstant bleibt. Zusätzlich zu lexikalischem und strukturellem Kasus gibt es semantischen Kasus. Dieser Kasus wird nicht von einem
regierenden Kopf wie einem Verb, Adjektiv oder einer Präposition zugewiesen, sondern ist auf eine bestimmte Funktion eines
Adverbials zurückzuführen. Beispielsweise stehen Zeitangaben wie \emph{den ganzen Tag} in (\mex{1})
im \ili{Deutsch}en im Akkusativ.
\ea
Er arbeitet den ganzen Tag.
\z
Da es in diesem Kapitel um das Passiv und seine Variation in den \ili{germanisch}en Sprachen geht, ignoriere ich hier den
semantischen Kasus.

\subsubsection{Nominative und Akkusativobjekte}

Bis jetzt wurde der Kasus, den ein Argument von seinem Kopf zugewiesen bekommt, in der Valenzliste des
Kopfes repräsentiert. Mit einer solchen Repräsentation bräuchten wir zwei verschiedene Lexikoneinträge für das Verb
\emph{lesen}: einen, in dem das Verb einen Nominativ und einen Akkusativ nimmt, wie in
(\mex{1}c), und einen, in dem es zwei Akkusative nimmt, wie in (\mex{1}d). (\mex{1}c) würde in
der Analyse von (\mex{1}a) verwendet und (\mex{1}d) in der Analyse von (\mex{1}b).
\eal
\ex 
Er        wird das        Buch lesen.
\ex 
Ich sah ihn das Buch lesen.
\ex \sliste{ NP[\type{nom}], NP[\type{acc}] }
\ex \sliste{ NP[\type{acc}], NP[\type{acc}] }
\zl
Statt zwei verschiedene, aber homophone Formen im Lexikon zu haben, kann man nur einen Lexikoneintrag vorschlagen und die eigentliche
Kasuszuweisung erst dann auflösen lassen, wenn der syntaktische Kontext genügend Information liefert. Je
nachdem, ob das Subjekt von \emph{lesen} als
Subjekt von \emph{wird} oder als Objekt von \emph{sah} realisiert wird, bekommt es also Nominativ oder Akkusativ. Solche
Kasus werden strukturelle Kasus genannt. Die Unterscheidung zwischen strukturellem und lexikalischem Kasus spielt eine
wichtige Rolle in der Analyse des Passivs. Es ist diese Unterscheidung, die eine einheitliche Analyse des
persönlichen und unpersönlichen Passivs möglich macht.

(\mex{1}) liefert weitere Beispiele und umfasst verschiedene Formen des Verbs (finit
vs.\ nicht-finit) sowie eine Nominalisierung:
\eal
\ex 
Der Installateur kommt.
\ex 
Der Mann lässt den Installateur kommen.
\ex 
das Kommen des Installateurs
\zl
Das Beispiel in (\mex{0}c) zeigt auch, dass das Subjekt von \emph{kommen} als
Genitiv realisiert werden kann. Somit sind Nominativ, Genitiv und Akkusativ strukturelle Kasus im \ili{Deutsch}en. (Die Frage,
ob einige oder alle Dative als struktureller Kasus behandelt werden sollten, wird weiter unten in Abschnitt~\ref{sec-dative-objects-tructurl-lexical} behandelt.)

Die Beispiele in (\mex{0}) zeigen, dass sich der Kasus von Subjekten im \ili{Deutsch}en ändern kann; die in (\mex{1})
zeigen, dass sich auch der Kasus von Akkusativobjekten ändern kann:
\eal
\ex 
Judit schlägt den Weltmeister.
\ex 
Der        Weltmeister    wird   geschlagen.
\zl

\subsubsection{Genitivobjekte}

Die Beispiele in (\mex{1}) zeigen Fälle von lexikalischem Kasus: Genitiv, der vom Verb abhängt, ist
lexikalisch, da er sich nicht ändert, wenn das Verb passiviert wird.
\eal
\ex[]{
Wir gedenken der Opfer.
}
\ex[]{
Der        Opfer   wird   gedacht.
}
\ex[*]{
Die Opfer wird / werden gedacht.
}
\zl
Wie das Beispiel in (\mex{0}c) zeigt, ist der Nominativ unmöglich. Das Genitivobjekt bleibt in
Passivkonstruktionen im Genitiv. Wie in Abschnitt~\ref{sec-impersonal-passive-phen} erläutert wurde,
werden Passive ohne Subjekt wie in (\mex{0}b) traditionell ``unpersönliche Passive'' genannt.\is{Passiv!unpersönliches}

\subsubsection{Dativobjekte}
\label{sec-dative-objects-tructurl-lexical}

Wenden\is{Kasus!Dativ|(} wir uns nun dem Dativ zu. Wenn wir Beispiele wie (\mex{1}) betrachten, sehen wir, dass sich der Dativ
im Passiv ebenfalls nicht ändert:
\eal
\ex 
Der Mann hat ihm geholfen.
\ex 
Ihm        wird geholfen.
\zl
In Analogie zu den obigen Genitivbeispielen sollte der Dativ also ein lexikalischer Kasus sein.

Es gibt aber Beispiele wie die in (\mex{1}), und gemäß der Auffassung, dass strukturelle Kasus
diejenigen Kasus sind, die je nach syntaktischer Umgebung variieren, sollte der Dativ ein struktureller Kasus sein.
\eal
\ex 
Der Mann  hat   den Ball dem Jungen geschenkt.
\ex 
Der Junge bekam den Ball geschenkt.
\zl
%\largerpage
Die Frage, ob der Dativ als struktureller oder lexikalischer Kasus anzusehen ist, wird heftig
diskutiert. Im Prinzip gibt es drei Möglichkeiten, und alle drei wurden in der
Literatur vorgeschlagen. Man könnte annehmen, dass alle Dative lexikalisch sind
\parencites[\page 80]{Haider85b}%
[\page 9]{Haider86}[\page 207, 217, 228]{HM94a}%
{Mueller99a,Mueller2001a}[\page 289]{Mueller2003e}%
[\page 97]{Scherpenisse86a}%
[\page 277, 291]{Pollard94a}%
%[]{Meurers99b} We are agnostic as to whether a complete structural case theory should also include
%the dative.
[67]{VS98a}[]{Abraham95a-u}%
[\page 187]{McIntyre2006a}% 
[Section~3]{Woolford2006a},\todostefan{wo Abraham?}
%-----------------------------------------------------------------------
dass einige lexikalisch und andere strukturell sind
\parencites
{Wegener85a}%
{Wegener90}% ganzes Papier 74--75, 79 100 Abschnitt 10 zu inheränten Dativen
% Im Japanischen wird angehobenes Subjekt zum Dativ, wenn Verbalkomplexbildung
                      % zu NP, NP, NP führt.
[\page 26]{denBesten85}% verbs like helfen, which assign dative case
[page 55--56]{denBesten85b}% Für helfen sagt er in Fn, dass die auch mit V' kombiniert werden.
[\page 161]{Fanselow87a}[\page 178, 205]{Fanselow2000b}[\page 181, 182, 206]{Fanselow2003b}% 181 lexical datives
[\page 286--287]{Czepluch88}%
[\page 77, 80]{Sternefeld95a}% für bekommen-Passiv mit helfen ist es dann doch struktureller Kasus
                             % und Er wird geholfen ist immer noch ausgeschlossen, weil er das mit
                             % einem Dativ-pro macht.
[\page 102--103]{Stechow96a}%
[\page 48, 51]{Wunderlich97a}%
%[\page 107]{Wunderlich97c}% dative passive possible -> structural lexical datives are not mentioned
%                          % but follow implicitely
[\page 553]{Molnarfi98a}% inheränte sind in V inkorporiert
%-----------------------------------------------------------------------
, oder dass alle Dative strukturell sind
\parencites[\page 80]{Sternefeld95a}%
[\page 203, 205--206]{Ryu97a}% SUBJ finit = nom, COMPS und INTARG = acc, COMPS und -INTARG = dat
[\page 96--97]{Gunkel2003b}.% two subtypes for structural case = nom v acc und nom v dat
% nom v acc müsste aber nom v acc v gen sein. Kasuszuweisung müsste dann an Objekte acc v dat sein.
% Zwei Merkmale, die GF auszeichnen: EA, IA. Passiv macht IA zu EA (IA kann auch leer sein).
% EA ist Nominativ, alle anderen sind -nom (und nicht Genitiv)
% Das funktioniert nicht, wenn das Isländische auch Dativpassive hat.
%\todostefan{update references}
%Molnarfi98a} muss noch kopiert werden

Ich folge \citet[9]{Haider86} und behandle alle Dative als lexikalische Kasus. Unter dieser Annahme wird der
Kontrast in den Beispielen in (\mex{1}) unmittelbar erklärt \citep[\page 20]{Haider86}:
% Haider hat andere Beispiele
\eal
\ex[]{
Er streichelt den Hund.
}
\ex[]{
Der        Hund wird  gestreichelt.
}
\ex[]{
sein Streicheln des    Hundes
}

\ex[]{\label{bsp-er-hilft-den-kindern}\iw{helfen}
Er hilft den Kindern.
}
\ex[]{
Den        Kindern  wird  geholfen.
}
\ex[]{
das Helfen der Kinder
}\label{das-helfen-der-Kinder}
\ex[*]{
sein Helfen der Kinder
}\label{sein-helfen-der-Kinder}
\zl
Das Akkusativobjekt von \emph{streicheln} kann im Passiv als Nominativ realisiert werden,
es ist also eindeutig ein struktureller Kasus. Nominalisierungen erlauben, dass dieses Objekt im Genitiv realisiert wird, wie (\mex{0}c) zeigt. Bei Dativen funktioniert das jedoch nicht. Das Dativobjekt von \emph{helfen}
kann nicht im Genitiv realisiert werden. (\ref{das-helfen-der-Kinder}) ist möglich, aber nur
mit einer Lesart, in der die Kinder das Agens sind, das heißt, die Nominalisierung in
(\ref{das-helfen-der-Kinder}) entspricht (\mex{1}) und nicht (\ref{bsp-er-hilft-den-kindern}):
\ea
Die Kinder helfen jemandem.
\z
Wenn das Agens durch ein pränominales Possessivum wie in (\ref{sein-helfen-der-Kinder}) ausgedrückt wird, ist der
Genitiv oder Dativ \emph{der Kinder} ausgeschlossen.

Die einzige Möglichkeit, den Dativ überhaupt auszudrücken, ist pränominal:
\ea
das Den-Kindern-Helfen
\z

%\largerpage[2]
\noindent
Daher haben Autoren, die annehmen, dass alle Dative strukturell sind, ein Problem damit, die Unterschiede bei
unpersönlichen Passiven und Nominalisierungen zu erklären.\footnote{\label{fn-structural-dative}
  Dieses Problem kann gelöst werden, indem man feinkörnigere Unterscheidungen innerhalb der strukturellen Kasus annimmt
  \citep[\page 96]{Gunkel2003b} und/oder zusätzliche Merkmale, die Akkusativobjekte herausgreifen
  \citep[\page 208]{Ryu97a}. Gunkels Ansatz ist gleichbedeutend mit der Aussage, dass primäre Objekte nominativisch
  oder dativisch sein müssen und sekundäre Objekte nominativisch oder akkusativisch. Für die Kasuszuweisung in verbalen
  Umgebungen besagt er, dass Objekte nicht-nominativisch sein müssen (S.\,112). Er schließt jedoch die
  Genitive in Nominalisierungen nicht ein. Würde man diese einbeziehen, wäre das sekundäre Objekt mit
  Nominativ, Genitiv und Akkusativ kompatibel. Für Objekte Nicht-Nominativ zu fordern, wäre nicht
  ausreichend, da dies die Option offen ließe, sekundäre Objekte in verbalen Umgebungen als Genitive zu realisieren,
  was ungrammatisch ist. Gunkel müsste festlegen, dass Objekte
  im Dativ oder Akkusativ stehen müssen, was zu einer recht komplexen und unattraktiven
  Darstellung führen würde, die die Fakten an verschiedenen Stellen der Grammatik wiederholt. Ryu nimmt Merkmale an, die
  das Subjekt und das sekundäre Objekt herausgreifen, und deshalb kann er zwischen primärem
  und sekundärem Objekt in dem für die Kasuszuweisung zuständigen Prinzip unterscheiden (S.\,205--206). Zur
  detaillierten Kritik an \citew{Ryu97a} und \citew{Gunkel2003b} siehe \citew[Abschnitt~3.4]{Mueller2003e}
  bzw.\ \citew[Kapitel~14.3.1]{MuellerLehrbuch3}.
} Zusätzlich gibt es ein Problem mit zweistelligen Verben. Während
einige Verben den Dativ nehmen, nehmen andere den Akkusativ, obwohl es kaum einen semantischen
Unterschied oder einen anderen Grund gibt, der dafür verantwortlich gemacht werden könnte.
\eal
\ex 
Er hilft ihm.
\ex 
Er unterstützt ihn.
\zl
Die Tatsache, dass \emph{helfen} ein Dativobjekt nimmt, während \emph{unterstützen} einen Akkusativ nimmt, ist
einfach eine Idiosynkrasie des \ili{Deutsch}en, die Sprecher des \ili{Deutsch}en lernen müssen, wenn sie die
Sprache erwerben. Die Information im Lexikoneintrag für \emph{helfen} muss sich also von der für
\emph{unterstützen} unterscheiden. Einige Autoren erkennen diesen Unterschied an und nehmen an, dass der Dativ
zweistelliger Verben lexikalisch ist, während der Dativ ditransitiver Verben strukturell ist \citep[\page 48, 51]{Wunderlich97a}. Die Annahme ist,
dass Verben ihrem ersten Argument den Nominativ zuweisen, ihrem letzten Argument den Akkusativ, und
wenn es ein zusätzliches Argument gibt, das weder das erste noch das letzte ist, bekommt es Dativ. Die
Vorhersage, die solche gemischten Ansätze machen, ist, dass das Dativpassiv bei
ditransitiven Verben möglich, bei zweistelligen Verben aber unmöglich sein sollte, da der Dativ für erstere
Verben strukturell und für letztere lexikalisch ist. Die empirische Situation ist nicht so eindeutig, wie man es sich
wünschen könnte. Einige Autoren akzeptieren Beispiele wie (\mex{1}); andere lehnen sie ab.
\eal
\label{ex-helfen-with-dative-passive}
\ex 
Er kriegte von vielen geholfen / gratuliert / applaudiert.
\ex 
Man kriegt täglich gedankt.
\zl
%\largerpage[2]
Es gibt jedoch belegte Beispiele:\footnote{
  Diese Beispiele wurden zuerst in \citew[\page 134--135]{Mueller2002b} und \citew[\page 293]{MuellerLehrbuch1} diskutiert.
}
\eal
\label{ex-helfen-with-dative-passive-corpus}
\ex "`Da kriege ich geholfen."'\footnote{
Frankfurter Rundschau, 26.06.1998, S.\,7.%
}

\ex
% auch nach applaudiert geholfen + bekommen und kriegen gesucht 21.09.2003
Heute morgen bekam ich sogar schon gratuliert.\footnote{%
Brief von Irene G.\ an Ernst G.\ vom 10.04.1943, Feldpost-Archive mkb-fp-0270.}

\ex
%Branich IG-Vorsitzender Friedel Schönel meinte deshalb, 
"`Klärle"' hätte es wirklich mehr als verdient, auch mal zu einem "`unrunden"' Geburtstag gratuliert zu bekommen.\footnote{
Mannheimer Morgen, 28.07.1999, Lokales; "`Klärle"' feiert heute Geburtstag.%
}

\ex 
Mit dem alten Titel von Elvis Presley "`I can't help falling in love"' bekam Kassier Markus Reiß zum Geburtstag gratuliert, [\ldots]\footnote{
Mannheimer Morgen, 21.04.1999, Lokales; Motor des gesellschaftlichen Lebens.%
}
     mit dem alten Titel von Elvis Presley \hphantom{"`}I can't help falling in love bekam Kassier Markus Reiß zum Geburtstag gratuliert\\

\glt `Kassier Markus Reiß wurde mit dem alten Elvis-Presley-Titel ``I can't help falling in love with you'' zum Geburtstag gratuliert.'
\zl
Es scheint, dass die Verben \emph{kriegen}, \emph{erhalten} und \emph{bekommen} auf dem Weg sind, zu
Hilfsverben zu werden. Ihre Bedeutung verblasst immer mehr. Daher gibt es kaum noch Selektionsbeschränkungen für das untergeordnete Verb. Die einzige Voraussetzung dafür, dass das Dativpassiv angewendet werden kann, ist
natürlich, dass das eingebettete Verb einen Dativ regiert.

%\largerpage[2]
Wenn das Dativpassiv nun bei zweistelligen Verben wie \emph{helfen} möglich ist und \emph{helfen}
einen lexikalischen Dativ regieren muss (da sonst der Unterschied zwischen \emph{helfen} und
\emph{unterstützen} nicht erklärt werden könnte)\footnote{%
  Siehe aber Fußnote~\ref{fn-structural-dative}. Es könnten zusätzliche Merkmale oder Subtypen des
  strukturellen Kasus angenommen werden.},
dann folgt, dass das Dativpassiv in der Lage sein muss,
einen lexikalischen Dativ in einen strukturellen Kasus zu verwandeln (der in den obigen Beispielen als Nominativ realisiert wird). Das
bedeutet, dass man annehmen könnte, dass alle Dative lexikalisch sind, sogar die Dative ditransitiver
Verben. Das erklärt, warum diese Dative in Passiven wie (\mex{1}) nicht als Nominative oder Akkusative realisiert werden:
\eal
\ex[]{ 
dass er dem Jungen den Ball gegeben hat
}
\ex[]{ 
dass dem Jungen der Ball gegeben wurde
}
\ex[*]{ 
dass der Junge den Ball gegeben wurde
}
\ex[*]{ 
dass den Junge der Ball gegeben wurde
}
\zl
Sie bleiben einfach im Dativ. Die einzige Ausnahme ist das Dativpassiv, und dies muss als Ausnahme betrachtet werden.%
\is{Kasus!Dativ|)}%

%% \subsubsubsection{Non-canonical accusative objects}

%% I already showed that the accusative can be a structural case. However, there are some exceptional
%% cases like those of subjectless verbs that govern an accusative (\mex{1}a) and ditransitive
%% verbs like \emph{lehren} `to teach' that take two accusative objects rather than one accusative and a dative.
%% \eal
%% \ex 
%% \gll Ihn dürstet.\\
%%      he.\ACC{} thirsty.is\\
%% \glt `He is thirsty.'
%% \ex 
%% \gll Der Vater lehrte seinen Sohn einen neuen Tritt.\\
%%      the.\NOM{} father taught his.\ACC{} son a.\ACC{} new kick\\
%% \glt `The father taught his son a new kick.'
%% \zl
%% Verbs like \emph{lehren} are generally bad in the passive.

%% \subsubsubsection{Adjektivumgebungen}

%% 
%% \subsubsubsection{Lexikalischer Kasus in Adjektivumgebungen}

%% Kasus von Objekten von Adjektiven kann sich nicht ändern.\\
%% Adjektive können Genitiv und Dativ zuweisen:
%% \eal
%% \ex Ich war mir \emph{dessen} sicher.
%% \ex Sie ist \emph{ihm} treu.
%% \zl
%% 
%% Die Zuweisung von Akkusativ ist ebenfalls möglich:
%% \eal
%% \ex Das ist \emph{diesen Preis} nicht wert.
%% \ex Der Student ist \emph{das Leben} im Wohnheim nicht gewohnt.\iw{gewohnt}\footnote{
%%         \citep*[S.\,312]{HB72a}
%%       }
%% \ex Du bist mir \emph{eine Erkl"arung} schuldig.\footnote{
%%         \citep*[S.\,620]{HFM81}
%%       }
%% \zl
%% Akkusativ ist bei Adjektivkomplementen aber selten \citep{Haider85b}.
%% }

%% 
%% \subsubsubsection{Struktureller Kasus in Adjektivumgebungen}

%% Kasus der Subjekte von Adjektiven hängt von der syntaktischen
%% Umgebung ab \citep{Wunderlich84}:
%% \eal
%% \ex \emph{Der Mond} wurde kleiner.\iw{klein}
%% \ex Er sah\iw{sehen} \emph{den Mond} kleiner werden.
%% \zl

%% }

%\if 0

%% \subsubsubsection{Semantische Kasus}
%% \label{sec-sem-kasus}
%% \is{Kasus!semantischer|(}

%% \subsubsubsection{Semantische Kasus}

%% \begin{itemize}
%% \item NPen können auch als Adjunkte auf"|treten \citep{Haider85b}:
%% \eal
%% \ex Sie hörten \emph{den ganzen Tag} dieselbe Schallplatte.
%% \ex Laßt \emph{mir} den Hund in Ruhe!
%% \ex \emph{Eines Tages} erschien ein Fremder.
%% \zl

%% \item auch der Urteilsdativ (\emph{Dativ iudicantis})  \citep{Wegener85b}:

%% \eal
%% \ex Das ist \emph{mir} zu\iw{zu!Grad} schwer.
%% \ex Das ist \emph{dem Kind} zu langweilig / nicht interessant genug.\iw{genug!Grad}
%% \ex Du läufst \emph{der Oma} zu\iw{zu!Grad} schnell.
%% \ex Das Wasser ist \emph{dem Baby} warm genug.\iw{genug!Grad}
%% \zl
%% \end{itemize}

%% \subsubsubsection{Zuweisung semantischer Kasus durch das Verb?}

%% \begin{itemize}
%% \item
%% Haider: Zuweisung durch Verb in (\mex{1}) nicht sinnvoll:
%% \ea
%% Sie hörten \emph{den ganzen Tag} dieselbe Schallplatte.
%% \z
%% Zeitangaben kommen auch in adjektivischen und nominalen Umgebungen vor:
%% % zitiert Toman83
%% \eal
%% \ex die Ereignisse \emph{letzten Sommer}
%% \ex der Flirt \emph{vorigen Dienstag}
%% \ex die \emph{diesen Sommer} sehr günstige Witterung
%% \ex die \emph{diesen Sommer} sehr teuren Urlaubsreisen
%% \zl
%% NPen mit strukturellem Kasus müssen in Nominalumgebungen
%% Genitiv sein. $\to$\\
%% In (\mex{0}) keine Zuweisung von strukturellem Kasus.

%% \item
%% Die Kasus in (\mex{0}) werden nicht aufgrund ihres Vorkommens in einer bestimmten
%% syntaktischen Struktur zugewiesen,\\
%% sondern sind vielmehr durch die Bedeutung des Nomens bestimmt.
%% \end{itemize}

%% \subsubsubsection{Akkusativ und Genitiv}

%% %\citep*{ZMT85a} -> semantische Kasusmarkierung
%% Der freie Akkusativ kommt bei Maß"-angaben\is{Maßangaben} (Zeitdauer und Zeitpunkt)
%% vor (\mex{1}) und Genitiv bei Lokalangaben oder Zeitangaben (\mex{2}).
%% \eal
%% \ex Sie studierte \emph{den ganze Abend}.
%% \ex \emph{Nächsten Monat}\iw{Monat} kommen wir.
%% \zl
%% \eal
%% \ex Ein Mann kam \emph{des Weges}.\iw{Weg}
%% \ex \emph{Eines Tages}\iw{Tag} sah ich sie wieder.
%% \zl

%% \subsubsubsection{Kongruenzkasus}

%% 
%% \subsubsubsection{Kongruenzkasus}

%% \begin{itemize}
%% \item Zwei Akkusative?
%% \eal
%% \ex Er nannte \emph{ihn} \rot{einen Experten}.
%% \ex \emph{Er} wurde \rot{ein Experte} genannt.
%% \zl
%% 
%% \item Wären das zwei unabhängige Akkusative,\\
%%       würde sich bei Passivierung nur einer ändern.

%% 
%% \item Kasus von \emph{einen Experten} wird \emph{Kongruenzkasus} genannt.\\
%% Die prädikative Phrase \emph{einen Experten} stimmt mit dem
%% Element,\\ über das prädiziert wird, im Kasus überein. 
%% \end{itemize}
%% }

%% 
%% \subsubsubsection{Kongruenzkasus mit Präpositionen}

%% Ähnliche Effekte kann man mit den Präpositionen \emph{als} und \emph{wie}
%% beobachten.
%% \eal
%% \ex \emph{Er} gilt als \rot{großer Künstler}.\footnote{
%%         \citew[S.\,203--204]{Heringer73a}.
%%       }
%% \ex Man lässt \emph{ihn} als \rot{großen Künstler} gelten\iw{gelten als}.
%% \zl
%% 
%% \eal
%% \ex Ich sehe \emph{ihn} als \rot{meinen Freund} an.\iw{ansehen}\footnote{
%%         \citew*[S.\,154]{SS88a}.
%% }
%% \ex \emph{Er} wird als \rot{mein Freund} angesehen.
%% \zl
%% }

%% 
%% \subsubsubsection{Kongruenzkasus mit Adjunkten}

%% Wie bei den prädikativen Argumenten gibt es auch Kongruenzkasus bei Adjunkten:
%% \eal
%% \ex Sie verhielt\iw{verhalten} \emph{sich} wie \emph{ihr Vater}.
%% \ex Ich behandelte\iw{behandeln} \emph{ihn} wie \emph{meinen Bruder}.
%% \ex Ich half\iw{helfen} \emph{ihm} wie \emph{einem Freund}.
%% \ex Ich erinnerte\iw{erinnern} mich \emph{dessen} wie \emph{eines fernen Alptraums}.
%% \zl

%% }

%% 
%% \subsubsubsection{Prädikation = Kasuskongruenz?}

%% \begin{itemize}
%% \item Kongruieren prädikative Phrasen immer mit dem Element,\\
%%       über das sie prädizieren?
%% \item Dies würde sofort auch Beispiele wie das in (\mex{1}) erklären:
%% \ea
%% Er wird ein großer Linguist.
%% \z
%% 
%% \item In AcI"=Konstruktionen müßten beide NPen im Akkusativ stehen.\\
%% Das ist nicht der Fall:
%% \eal
%% %\ex Laß ihn einen großen Linguisten werden.\label{bsp-lass-ihn-einen-grossen}
%% \ex Laß\iw{lassen|(} den wüsten Kerl [\ldots] meinetwegen ihr Komplize sein.\footnote{
%%         (\ref{bsp-lass-den-wuesten-kerl}) und (\ref{bsp-lass-mich}) sind aus dem \citet*[{\S}\,6925]{Duden66}.\iaf{Duden} %\citet*[{\S}\,1473]{Duden73}.\iaf{Duden}
%%         Die Quellen finden sich dort.
%%       }\label{bsp-lass-den-wuesten-kerl}
%% \ex Laß mich dein treuer Herold sein.\label{bsp-lass-mich}
%% \ex Baby, laß\iw{lassen|)} mich dein Tanzpartner sein.\footnote{
%%         Funny van Dannen, Benno-Ohnesorg-Theater, Berlin, Volksbühne, 11.10.1995
%%         }
%% \zl
%% 
%% \item
%% $\to$ Nominativ des Nicht-Subjekts in Kopulakonstruktionen ist\\
%%       ein lexikalischer Kasus \citep[S.\,54]{Thiersch78a}.
%% \end {itemize}
%% }

%% \subsubsubsection{Der Kasus nicht ausgedrückter Subjekte}
%% \label{sec-kasus-nicht-realisierter-subj}

%% 
%% \subsubsubsection{Der Kasus nicht ausgedrückter Subjekte (I)}
%% \savespace

%% \begin{itemize}
%% \item \citet*[Kapitel~6]{Hoehle83}:\\
%% Kasus nicht an der Oberfläche auf"|tretender Elemente bestimmbar.

%% {\em ein- nach d- ander-\/} kann sich auf mehrzahlige Konstituenten beziehen. 

%% Dabei muß Kasus und Genus mit der Bezugsphrase übereinstimmen.
%% 
%% \item In (\mex{1}) Bezug auf Subjekte bzw.\ Objekte:
%% \eal
%% \ex Die Türen sind eine nach der anderen kaputtgegangen.
%% \ex Einer nach dem anderen haben wir die Burschen runtergeputzt.
%% \ex Einen nach dem anderen haben wir die Burschen runtergeputzt.
%% \ex Ich ließ die Burschen einen nach dem anderen einsteigen.
%% \ex Uns wurde einer nach der anderen der Stuhl vor die Tür gesetzt.
%% \zl
%% \end{itemize}
%% }

%% 
%% \subsubsubsection{Der Kasus nicht ausgedrückter Subjekte (II)}
%% \savespace

%% In (\mex{1}) Bezug auf Dativ- bzw.\ Akkusativobjekte
%% eingebetteter Infinitive:

%% \eal
%% \ex Er hat uns gedroht, die Burschen demnächst einen nach dem anderen wegzuschicken.
%% \ex Er hat angekündigt, uns dann einer nach der anderen den Stuhl vor die Tür zu setzen.
%% \ex Es ist nötig, die Fenster, sobald es geht, eins nach dem anderen auszutauschen.
%% \zl

%% }

%% 
%% \subsubsubsection{Der Kasus nicht ausgedrückter Subjekte (III)}
%% \savespace

%% In (\mex{1}) Bezug auf Subjekt innerhalb der Infinitiv"=VP:
%% \eal
%% \ex Ich habe den Burschen geraten, im Abstand von wenigen Tagen einer nach dem anderen
%%       zu kündigen.
%% \ex Die Türen sind viel zu wertvoll, um eine nach der anderen verheizt zu werden.
%% \ex Wir sind es leid, eine nach der anderen den Stuhl vor die Tür gesetzt zu kriegen.
%% \ex Es wäre fatal für die Sklavenjäger, unter Kannibalen zu fallen und einer nach dem
%%       anderen verspeist zu werden.
%% \zl
%% {\em ein- nach d- ander-\/} im Nominativ $\to$\\
%% Das nicht realisierte Subjekt steht ebenfalls im Nominativ.

%% }

%% 
%% \subsubsubsection{Der Kasus nicht ausgedrückter Subjekte (IV)}

%% Dasselbe gilt für nicht realisierte Subjekte von adjektivischen Partizipien:
%% \eal
%% \ex die eines nach dem anderen einschlafenden Kinder
%% \ex die einer nach dem anderen durchstartenden Halbstarken
%% \ex die eine nach der anderen loskichernden Frauen
%% \zl
%% }

%% 
%% \subsubsubsection{Der Kasus nicht ausgedrückter Subjekte (V)}

%% Man muß also sicherstellen, daß auch nicht realisierte Subjekte Kasus zugewiesen bekommen.
%% Würde man diesen Kasus unterspezifiziert lassen, würden Sätze wie (\mex{1}) falsch analysiert werden.
%% \judgewidth{\#}
%% \ea[\#]{
%% Ich habe den Burschen geraten, im Abstand von wenigen Tagen einen nach dem anderen zu kündigen.
%% }
%% \z
%% In der zulässigen Lesart von (\mex{0}) ist die Phrase 
%% \emph{einen nach dem anderen} das Objekt von \emph{kündigen} und kann
%% sich nicht auf das Subjekt des Infinitivs, das referenzidentisch
%% mit \emph{den Burschen} ist, beziehen.

%% }

Nach dieser Diskussion des lexikalischen und strukturellen Kasus im \ili{Deutsch}en wenden wir uns nun dem Kasusprinzip zu,
das für die Kasuszuweisung zuständig ist. Wie in Abschnitt~\ref{sec-linking} erläutert wurde, wird angenommen, dass alle Argumente eines
Kopfes in einer Liste repräsentiert werden: der \textsc{argument structure}-Liste (\argst-Liste\isfeat{arg-st}).
(\mex{1}) zeigt die Argumentstrukturliste eines ditransitiven Verbs wie \word{geben}:
\ea
\sliste{ NP[\type{str}], NP[\type{ldat}], NP[\type{str}] }
\z
Wie oben argumentiert wurde, wird der Dativ als lexikalischer Kasus behandelt. \type{ldat} ist eine Abkürzung für lexikalischen
Dativ und \type{str} steht für strukturellen Kasus. Das Kasusprinzip
hat die folgende Form (angepasst aus \citealp{Prze99}; \citealp{Meurers99b}):

\begin{principle-break}[\hypertarget{case-p}{Kasusprinzip}]\is{Prinzip!Kasus}
\label{case-p}
\begin{itemize}
\item In einer Liste, die sowohl das Subjekt als auch die Komplemente eines verbalen Kopfes enthält, bekommt das erste
  Element mit strukturellem Kasus den Nominativ, sofern es nicht von einem höheren Kopf angehoben wird.
\item Alle anderen Elemente in dieser Liste, die strukturellen Kasus haben und nicht angehoben werden, bekommen Akkusativ.\is{Kasus!Akkusativ}
\item In nominalen Umgebungen bekommen Elemente mit strukturellem Kasus Genitiv.\is{Kasus!Genitiv}
\end{itemize}
\end{principle-break}
Dieses Prinzip ist von \citet*{YMJ87} inspiriert und funktioniert, wie weiter unten gezeigt wird, für alle
hier analysierten Sprachen, insbesondere für das komplexe Kasussystem des \ili{Isländisch}en. Das hier angenommene Kasussystem
unterscheidet sich darin, dass es Elementen, die zu einem höheren Prädikat angehoben werden, keinen Kasus zuweist. Dieser Punkt wird weiter unten genauer erläutert.

\largerpage
Der Effekt dieses Prinzips wird anhand der Verben in (\mex{1}) erläutert:
\ea
\begin{tabular}[t]{@{}l@{~}l@{~}l@{}}
a. & \emph{schläft}:       & \argst \sliste{ NP[\type{str}]$_i$ }\\[2mm]
b. & \emph{unterstützt}: & \argst \sliste{ NP[\type{str}]$_i$, NP[\type{str}]$_j$ }\\[2mm]
c. & \emph{hilft}:          & \argst \sliste{ NP[\type{str}]$_i$, NP[\type{ldat}]$_j$ }\\[2mm]
d. & \emph{schenkt}:     & \argst \sliste{ NP[\type{str}]$_i$, NP[\type{ldat}]$_j$, NP[\type{str}]$_k$ }\\
\end{tabular}
\z
Das erste Element in diesen Listen, das strukturellen Kasus hat, bekommt Nominativ und das zweite
Akkusativ. Das ist genau das, was man erwartet. Das Ergebnis ist in (\mex{1}) angegeben. \type{snom} steht
für strukturellen Nominativ und \type{sacc} für strukturellen Akkusativ.
\ea
%\oneline{%
\scalebox{.99}{%
\begin{tabular}[t]{@{}l@{~}l@{~}l@{}}
a. & \emph{schläft}:       & \argst \sliste{ NP[\type{snom}]$_i$ }\\[2mm]
b. & \emph{unterstützt}: & \argst \sliste{ NP[\type{snom}]$_i$, NP[\type{sacc}]$_j$ }\\[2mm]
c. & \emph{hilft}:          & \argst \sliste{ NP[\type{snom}]$_i$, NP[\type{ldat}]$_j$ }\\[2mm]
d. & \emph{schenkt}:     & \argst \sliste{ NP[\type{snom}]$_i$, NP[\type{ldat}]$_j$, NP[\type{sacc}]$_k$ }\\
\end{tabular}
}
\z
\is{Kasus!struktureller|)}\is{Kasus!lexikalischer|)}

\subsection{Argumentreduktion und Kasuszuweisung: Das Passiv}
\label{sec-case-assignment-passive}

Angesichts der Unterscheidung zwischen strukturellem und lexikalischem Kasus ist die Analyse des Passivs wirklich einfach und
entspricht direkt der Intuition, dass das Passiv die Unterdrückung des Subjekts ist (des prominentesten, also des ersten Arguments in der \argstl).\footnote{
  Einige Autoren nehmen eine Degradierung des Subjekts an, das heißt, das Subjekt wird in ein Komplement umgewandelt, eine \emph{by}-PP
  oder eine \emph{von}-PP \citep[\page 216]{ps}. Ich folge \citet[\page161]{Hoehle78a}, \citet{Sadzinski87a}, \citet[\page174]{Stechow90a},
\citet[\page255]{Zifonun92a}, 
%\NOTE{\citet[\page86]{Leiss92a},}
\citet[\page181]{Lieb92a},
\citet[\page740]{Wunderlich93a}, 
%\NOTE{\citet[\page150]{Primus99a},}
%cannot follow myself \citet{Mueller2003e}, 
\citet[\page 65]{Gunkel2003b} und anderen darin, dass die PPs zum Ausdruck des Agens Adjunkte sind. Siehe
  \citew[Abschnitt~5]{Mueller2003e} für Details.
} Wenn das erste Argument aus den
Listen in (\mex{-1}) entfernt wird, ergeben sich die folgenden Listen:
\ea
\begin{tabular}[t]{@{}l@{~}l@{~}l}
a. & \emph{geschlafen}:  & \argst \sliste{ }\\[2mm]
b. & \emph{unterstützt}: & \argst \sliste{ NP[\type{str}]$_j$ }\\[2mm]
c. & \emph{geholfen}:    & \argst \sliste{ NP[\type{ldat}]$_j$ }\\[2mm]
d. & \emph{geschenkt}:   & \argst \sliste{ NP[\type{ldat}]$_j$, NP[\type{str}]$_k$ }\\
\end{tabular}
\z
Die NPs, die in (\mex{0}) an erster Position stehen, standen in (\mex{-1}) an zweiter Position. Die
erste NP mit strukturellem Kasus bekommt Nominativ, und somit ergeben sich die folgenden Kasuszuweisungen:
\ea
\begin{tabular}[t]{@{}l@{~}l@{~}l}
a. & \emph{geschlafen}:  & \argst \sliste{ }\\[2mm]
b. & \emph{unterstützt}: & \argst \sliste{ NP[\type{snom}]$_j$ }\\[2mm]
c. & \emph{geholfen}:    & \argst \sliste{ NP[\type{ldat}]$_j$ }\\[2mm]
d. & \emph{geschenkt}:   & \argst \sliste{ NP[\type{ldat}]$_j$, NP[\type{snom}]$_k$ }\\
\end{tabular}
\z
Da lexikalischer Kasus wie in (\mex{0}c--d) vom Kasusprinzip nicht betroffen ist, bleibt er so, wie er
spezifiziert wurde, nämlich Dativ.

\largerpage
Es sei hier angemerkt, dass dieser einfache Ansatz zum Passiv sowohl das sogenannte
persönliche als auch das unpersönliche Passiv erfasst. Die Passive von \emph{schlafen} und \emph{helfen} werden unpersönliche Passive genannt, da die jeweiligen Sätze kein Subjekt haben. 
\eal
\ex 
dass geschlafen wurde
\ex
dass dem Mann geholfen wurde
\zl
Die Passive von \emph{unterstützen} und \emph{schenken} haben hingegen
Subjekte, nämlich die Argumente, die im Aktiv als Akkusativobjekte realisiert werden:
\eal
\ex 
dass der Mann unterstützt wurde
\ex
dass dem Jungen der Ball geschenkt wurde
\zl
Diejenigen Analysen, die alle Kasus lexikalisch zuweisen, müssten annehmen, dass der Kasus der Objekte
(Akkusativ) im Passiv in Nominativ geändert wird. Es gäbe also zwei Varianten des
Passivs: Das unpersönliche Passiv unterdrückt nur das Subjekt, und das persönliche Passiv unterdrückt das
Subjekt und ändert zusätzlich den Kasus des Objekts in Nominativ. Die Analyse, die die
Unterscheidung zwischen strukturellem und lexikalischem Kasus verwendet, verschiebt die Kasuszuweisung lediglich bis zu dem Punkt, an dem
klar ist, was der richtige Kasus sein wird. Wenn wir ein Partizip haben und es mit dem Passivhilfsverb verwenden, ist
klar, welchen Kasus die Argumente haben müssen.

%% \subsubsubsubsection{Dativpassiv}

%% \frame[shrink=15]{
%% \subsubsection{Dativpassiv}

%% Bei der Kombination von \emph{geholfen} und
%% \emph{bekommen} bzw.\ von \emph{geschenkt} und \emph{bekommen} wird das Dativargument von 
%% \emph{geholfen} bzw.\ von \emph{geschenkt} zum ersten Argument gemacht und der lexikalische
%% Dativ beim eingebetteten Verb wird zu einem strukturellen Kasus beim Passiv"=Hilfsverb:
%% \ea
%% \begin{tabular}[t]{@{}l@{~}l@{~}l}
%% c. & \emph{hilft}:       & \argst \sliste{ NP[\type{str}]$_j$, NP[\type{ldat}]$_k$ }\\
%% d. & \emph{schenkt}:     & \argst \sliste{ NP[\type{str}]$_j$, NP[\type{str}]$_k$, NP[\type{ldat}]$_l$ }\\
%% \end{tabular}
%% \z
%% \ea
%% \begin{tabular}[t]{@{}l@{~}l@{~}l}
%% a. & \emph{geholfen bekommt}:    & \argst \sliste{ NP[\type{str}]$_k$ }\\
%% b. & \emph{geschenkt bekommt}:   & \argst \sliste{ NP[\type{str}]$_l$, NP[\type{str}]$_k$ }\\
%% \end{tabular}
%% \z
%% Details kommen im Kapitel über Passiv.

%% Kasusvergabe: Dadurch, daß das Dativargument an erster Stelle in der Valenzliste\\
%% von \emph{geholfen bekommen} bzw.\ von \emph{geschenkt bekommen} steht, kriegt es
%% Nominativ. 

%% Bei \emph{geschenkt bekommen} bekommt das zweite Element (das direkte Objekt) Akkusativ.

%% Die Umwandlung eines lexikalischen in einen strukturellen Kasus ist unschön,\\
%% es scheint zur Zeit jedoch keine bessere Alternative zu geben. 

%% }

\subsection{Argumenterweiterung und Kasuszuweisung: AcI-Konstruktionen}

Das Kasusprinzip enthält Beschränkungen für die Kasuszuweisung, die die Zuweisung an
angehobene Elemente verbieten. Diese Beschränkungen wurden noch nicht erläutert. Betrachten wir die Beispiele in (\mex{1}):
\eal
\ex
Der Junge liest den Aufsatz.
\ex\label{ex-the-man-lets-the-boy-read-the-book}
Der Mann lässt den Jungen den Aufsatz lesen.
\zl
Das Beispiel (\mex{0}a) zeigt, dass dem Subjekt von \emph{lesen}
Nominativ zugewiesen wird. Das Subjekt von \emph{lesen} bekommt in (\mex{0}b) jedoch Akkusativ. Wenn
man also den Kasus auf Basis der Argumentstruktur von \emph{lesen} in (\mex{0}b) zuweisen würde, würde man
Nominativ zuweisen, aber das AcI-Verb\is{Verb!AcI}\is{Verb!kausatives} \emph{lassen} weist seinem Objekt Akkusativ zu. Der
Punkt ist, dass das Subjekt von \emph{lesen} zum Objekt von \emph{lassen} angehoben wird. Das Kasusprinzip
ist so eingerichtet, dass Kasus nur denjenigen Argumenten zugewiesen wird, die nicht zu einem höheren Kopf angehoben werden. Daher bekommt \emph{den Jungen} seinen Kasus nicht von \emph{lesen}, sondern von \emph{lässt}.

Die Analyse von (\mex{0}b) ist in Abbildung~\vref{fig-vc-aci} angegeben.
\begin{figure}
\centerfit{
\begin{forest}
sm edges
[V\feattab{
              \sliste{ }}
        [{NP[\type{snom}]} [der Mann, roof] ]
        [V\feattab{
              \sliste{ NP[\type{snom}] }}
          [{NP[\type{sacc}]} [den Jungen, roof] ]
          [V\feattab{
%              \vform \type{fin},\\
              \sliste{ NP[\type{snom}], NP[\type{sacc}]}} 
            [{NP[\type{sacc}]} [den Aufsatz, roof] ]
            [V\feattab{
%                \vform \type{fin},\\
                \sliste{ NP[\type{snom}], NP[\type{sacc}], NP[\type{sacc}]}} 
               [~~~~~~V\feattab{
%                \vform \type{bse},\\
                 \sliste{ NP[\type{sacc}], NP[\type{sacc}]}} [lesen] ]
               [V\feattab{
%                \vform \type{fin},\\
                  \sliste{ NP[\type{snom}], NP[\type{sacc}], NP[\type{sacc}], V }} [lässt] ] ] ] ] ]
\end{forest}}
\caption{\label{fig-vc-aci}Analyse von AcI-Konstruktionen als Anhebungskonstruktionen und der Verbalkomplex im Deutschen}
\end{figure}
Die Argumente von \emph{lesen} werden von \emph{lässt} übernommen. Da \emph{lässt}
ein eigenes Argument beisteuert, den Verursacher bzw.\ denjenigen, der die Erlaubnis gibt, selegiert \emph{lässt}
in dem konkreten in Abbildung~\ref{fig-vc-aci} dargestellten Satz drei NPs mit strukturellem Kasus und ein Verb. Gemäß dem Kasusprinzip bekommt die erste NP mit strukturellem Kasus
Nominativ und die anderen NPs mit strukturellem Kasus bekommen Akkusativ. Das ergibt eine Liste mit einer
NP im Nominativ und zwei NPs im Akkusativ.

(\mex{1}) zeigt die \argstl von \emph{lässt}, wenn es mit \emph{schlafen},
\emph{unterstützen}, \emph{helfen} bzw.\ \emph{schenken} kombiniert wird.
\ea
\onelineea{%
\begin{tabular}[t]{@{}l@{~}l@{~}l@{}}
a. & \emph{lässt} + \emph{schlafen}:      & \argst \sliste{ NP[\type{str}]$_l$, NP[\type{str}]$_i$, V}\\[2mm]
b. & \emph{lässt} + \emph{unterstützen}: & \argst \sliste{ NP[\type{str}]$_l$, NP[\type{str}]$_i$, NP[\type{str}]$_j$, V }\\[2mm]
c. & \emph{lässt} + \emph{helfen}:        & \argst \sliste{ NP[\type{str}]$_l$, NP[\type{str}]$_i$, NP[\type{ldat}]$_j$, V }\\[2mm]
d. & \emph{lässt} + \emph{schenken}:      & \argst \sliste{ NP[\type{str}]$_l$, NP[\type{str}]$_i$, NP[\type{ldat}]$_j$, NP[\type{str}]$_k$, V }\\
\end{tabular}
}
\z
\largerpage[-1]
Die hinzugefügte NP hat den Index \type{l}. Als erste NP mit strukturellem Kasus auf diesen Listen
bekommt sie Nominativ. Alle anderen Elemente dieser Liste mit strukturellem Kasus bekommen Akkusativ. Daher wird dem
Subjekt des eingebetteten Verbs Akkusativ zugewiesen, die lexikalischen Kasus bleiben gleich und die
Akkusativobjekte des eingebetteten Verbs bekommen ebenfalls Akkusativ, da ihr Kasus auch strukturell ist.

%\largerpage
Man beachte, dass die Frage, ob eine Sprache einen Verbalkomplex hat oder nicht, orthogonal zu Fragen der
Kasuszuweisung ist. Abbildung~\vref{fig-the-man-lets-the-boy-read-the-book} zeigt die Analyse der \ili{englisch}en Übersetzung von (\ref{ex-the-man-lets-the-boy-read-the-book}).
\begin{figure}
\centerfit{
\begin{forest}
sm edges
[S
        [{NP[\type{snom}]} [the man, roof] ]
        [VP
          [\vbar [V [lets] ]
                 [{NP[\type{sacc}]} [the boy,roof] ] ]
              [VP
                [V [read] ]
                [{NP[\type{sacc}]} [the book, roof] ] ] ] ]
\end{forest}}
\caption{\label{fig-the-man-lets-the-boy-read-the-book}AcI-Konstruktionen im Englischen}
\end{figure}
\emph{let} selegiert das Subjekt, das Objekt und eine VP. Das Subjekt von \emph{read} ist
gleichzeitig das Objekt von \emph{let}, und daher weist das Kasusprinzip nicht dem
Subjekt des eingebetteten Verbs \emph{read} Nominativ zu, sondern dem Objekt des Matrixverbs \emph{let} Akkusativ.%
\is{Kasus|)}

%% 
%% \subsubsubsection{Adjektivsubjekte}

%% Die Kasuszuweisungen an das Subjekt von Adjektiven funktioniert analog. Die Kopula wird mit dem Adjektiv
%% verbunden, und es entsteht eine Valenzliste, die die Argumente des Adjektivs enthält (\mex{1}a).\\
%% Wird ein solcher Komplex noch unter ein AcI"=Verb wie \emph{sehen} eingebettet,\\
%% erhält man die Liste in (\mex{1}b):
%% \ea
%% \begin{tabular}[t]{@{}l@{~}l@{~}l}
%% a. & \emph{kleiner werden}:     & \argst \sliste{ NP[\type{str}]$_j$ }\\
%% b. & \emph{kleiner werden sah}: & \argst \sliste{ NP[\type{str}]$_i$, NP[\type{str}]$_j$ }\\
%% \end{tabular}
%% \z
%% Die Kasuszuweisung funktioniert analog zu den bereits diskutierten Fällen. In den verbalen Umgebungen
%% der Kopula bzw.\ des AcI"=Verbs bekommen die NPen mit strukturellem Kasus Nominativ bzw.\ Akkusativ.%
%% }

%% \subsubsubsection{Semantischer Kasus}
%% 
%% \subsubsubsection{Semantischer Kasus (I)}

%% Der Kasus von NPen wie \emph{den ganzen Tag} in (\mex{1}) ist von der syntaktischen Umgebung unabhängig.
%% \eal
%% \ex Sie arbeiten den ganzen Tag.
%% \ex Den ganzen Tag wird gearbeitet, [\ldots].\footnote{
%%   \url{http://www.philo-forum.de/philoforum/viewtopic.html?p=146060}. \urlchecked{12}{05}{2005}.
%% }
%% \zl
%% Daß die NP im Akkusativ steht, hängt mit ihrer Funktion zusammen. 

%% Unterschiedliche Lexikoneinträge für \emph{Tag} in (\mex{0}) und (\mex{1}):
%% \eal
%% \ex Ich liebe diesen Tag.
%% \ex Dieser Tag gefällt mir.
%% \zl

%% In (\mex{0}) liegen ganz gewöhnliche Argumente vor,\\
%% in (\mex{-1}) dagegen ein Adjunkt. 
%% }

%% 
%% \subsubsubsection{Semantischer Kasus (II)}

%% Adjunkte unterscheiden sich von Argumenten durch ihren \modw und durch ihern \contw.

%% Für (\mex{-1}) muß es unter \cont eine Dauer-Relation geben.

%% Zusammen mit dieser Information wird im Lexikoneintrag für das modifizierende Nomen der Kasus fest kodiert. 

%% Die morphologische Komponente kann dann für diesen Lexikoneintrag nur die Akkusativform erzeugen, 
%% da alle anderen Flexionsformen mit der bereits im Lexikoneintrag angegebenen Kasusinformation inkompatibel sind. 

%% Dadurch wird sichergestellt, daß Sätze wie (\mex{1}) nicht analysiert werden:
%% \eal
%% \ex[*]{
%% Er arbeitet der ganze Tag.
%% }
%% \ex[*]{
%% weil der ganze Tag gearbeitet wurde
%% }
%% \zl

%% }

\subsection{Erklärung der einzelsprachübergreifenden Unterschiede}

Da die Argumentreduktion und Kasuszuweisung für das \ili{Deutsch}e bereits in
Abschnitt~\ref{sec-case-assignment-passive} erläutert wurde, möchte ich nun direkter darauf eingehen und
Lexikoneinträge für das Passiv- und Perfekthilfsverb des \ili{Deutsch}en angeben. Danach bespreche ich die anderen
Sprachen (\zb das \ili{Dänisch}e, \ili{Englisch}e und \ili{Isländisch}e) und erkläre, wie sich die Unterschiede analytisch erfassen lassen.

\subsubsection{Reduktion des designierten Arguments}

\largerpage
\citet[\page 10]{Haider86} hat vorgeschlagen, das Argument eines Verbs zu markieren, das Subjekteigenschaften hat. Er nennt
diese besonderen Argumente \emph{designiertes Argument}. \citet{HM94a} haben diese Idee auf HPSG übertragen, und
\citet{Mueller2003e} hat sie leicht modifiziert, um bestimmte Fakten bei Modalinfinitiven
richtig zu erfassen. Eine wichtige Verwendung des designierten Arguments besteht darin, so"=genannte
unakkusativische Verben\is{Verb!unakkusativisches} von unergativischen\is{Verb!unergativisches} Verben zu unterscheiden. \citet{Perlmutter78} hat darauf hingewiesen, dass unakkusativische Verben
bemerkenswerte Eigenschaften haben, und argumentiert, dass ihre Subjekte nicht wirklich Subjekte sind, sondern sich eher
wie Objekte verhalten. Eine ihrer Eigenschaften ist, dass sie keine Passive erlauben. Außerdem können ihre
Partizipien attributiv verwendet werden, was bei unergativischen Verben nicht möglich ist:

\eal
\ex[]{ 
der angekommene Zug
}
\ex[*]{
der geschlafene Mann
}
\zl

\noindent
Dies wird erklärt, wenn man annimmt, dass das Subjekt von \emph{ankommen} objektartige
Eigenschaften hat und sich daher wie das Objekt transitiver Verben verhält:
\ea
der geliebte Mann
\z
%\largerpage
\emph{Mann} füllt die Objektstelle von \emph{geliebte}. Wenn das einzige Argument von \emph{ankommen}
als Objekt behandelt wird, ist die Ähnlichkeit zum transitiven \emph{lieben} sofort
erklärt. Ebenso wird die Tatsache erklärt, dass Unakkusativa\is{Verb!unakkusativisches} keine Passive erlauben: Wenn
das Passiv die Unterdrückung des \isi{Subjekt}s ist und \emph{ankommen} in diesem Sinne kein Subjekt hat, kann das
Passiv nicht angewendet werden.

\eal
\ex[]{
Der Zug ist angekommen.
}
\ex[*]{
weil angekommen wurde
}
\zl

\noindent
In den HPSG-Analysen nehmen die Autoren an, dass es ein listen"=wertiges Merkmal \textsc{designated
  argument} (\textsc{da}) gibt. Diese Liste enthält das Subjekt transitiver und unergativischer Verben
(intransitiver Verben, die nicht unakkusativisch sind). Der \dav unakkusativischer\is{Verb!unakkusativisches} Verben ist die leere Liste,
da diese Verben kein Argument mit Subjekteigenschaften haben.

Das Passiv wird als Lexikonregel analysiert, die einen Lexikoneintrag für das Partizip lizenziert. Die
\argstl des Partizips ist die \argstl des Verbstamms, der die Eingabe der Lexikonregel ist, minus
die \dalist. Da dies nicht der Fokus dieses Buches ist, werde ich unakkusativische Verben im
Folgenden nicht diskutieren. (\mex{1}) liefert einige prototypische Beispiele für unergativische und transitive Verben:

%\largerpage[-1]
\ea
%\resizebox{\linewidth}{!}{%
\begin{tabular}[t]{@{}l@{ }l@{ }l@{ }l@{}}
  &                     & \textsc{arg-st} & \textsc{da}\\[2mm]
a.&tanzen:   & \sliste{ \ibox{1}NP[\type{str}] }                                              & \sliste{ \ibox{1} }\\[2mm]
b.&lesen:    & \sliste{ \ibox{1}NP[\type{str}], NP[\type{str}] }                              & \sliste{ \ibox{1} }\\[2mm]
c.&schenken: & \sliste{ \ibox{1}NP[\type{str}], NP[\type{ldat}], NP[\type{str}] } & \sliste{ \ibox{1} }\\[2mm]
d.&helfen:   & \sliste{ \ibox{1}NP[\type{str}], NP[\type{ldat}] }                            & \sliste{ \ibox{1} }\\
\end{tabular}
%}
\z
Die Lexikonregel\is{$\mapsto$}, die das Partizip bildet, ist in (\ref{lr-passive-prelim}) skizziert:

\ea
\label{lr-passive-prelim}
Lexikonregel\is{Lexikonregel!Passiv} zur Bildung des Partizips (vorläufig):\\
\ms[stem]{
head   & \ms[verb]{ da & \ibox{1}\\
                  }\\
arg-st & \ibox{1} $\oplus$ \ibox{2} \\
} $\mapsto$
\ms[word]{
arg-st & \ibox{2} \\
}
\z
Diese Regel teilt die \argstl der Eingabe in zwei Listen \ibox{1} und \ibox{2} auf. \ibox{1} ist
identisch mit dem \dav. Daher wird das designierte Argument von der \argstl entfernt und ist nicht in
dem von der Regel lizenzierten Lexikoneintrag vorhanden. Die Lexikonregel bildet einen \type{stem} auf ein
\type{word} ab. \stem{geb} wird beispielsweise auf das Partizip \emph{gegeben} abgebildet. Um die
Dinge einfach zu halten, wird die Abbildung der phonologischen Information in der Regel in
(\mex{0}) nicht gezeigt. %I will provide phonological information under \phon in the figures below.
\label{page-lexical-rule-explanation-start}%
Die Lexikonregel erzeugt neue Lexikoneinträge. Sie wird auf Stämme angewendet, die im Lexikon aufgeführt sind, und auf
Stämme, die von anderen Lexikonregeln lizenziert werden. \stem{geb} und \stem{säg} sind beispielsweise im Lexikon aufgeführt. Der Stamm für \stem{zersäg} wird aus \stem{säg} durch eine andere
Lexikonregel abgeleitet, die das Präfix \prefix{zer} hinzufügt. Dieser Stamm kann Eingabe der Passivierungs-Lexikonregel sein,
die dann die Partizipform \emph{zersägt} ausgibt, wie sie in (\mex{1}c) verwendet wird.
\eal
\ex 
Aicke sägt.
\ex
Aicke zersägt den Baum.
\ex
Der Baum wurde zersägt.
\zl
Lexikonregeln, wie sie hier angenommen werden, sind vollständig in HPSG integriert und können als unär
verzweigende Bäume dargestellt werden. Für die technischen Details siehe \citet{Meurers2000b} und \citet{BC99a}. Man beachte, dass jeder Knoten in den in diesem Buch angegebenen Bäumen seine eigene
phonologische Information trägt. Dies weicht von normalen Phrasenstrukturgrammatiken ab, in denen nur die
phonologische Information ganz unten in den Bäumen (an den Blättern) relevant ist und der gesamte Baum
für eine Verkettung der gesamten phonologischen Information der Blätter steht. Abbildung~\ref{fig-LR-as-tree}
zeigt die Baumdarstellung für die Abbildung vom Stamm \stem{säg} über \stem{zersäg} zum
Passivpartizip \emph{zersägt}.\footnote{%
Mehr zur Morphologie\is{Morphologie} innerhalb von HPSG siehe \citew[Kapitel~19]{MuellerLehrbuch}.
} Strukturteilungen werden weggelassen, um die Lesbarkeit zu wahren.

\begin{figure}
%\vfill
% \begin{forest}
% %sm edges
% [{\ms[word]{\phon  & \phonliste{ given }\\[2mm]
%             \spr   & \sliste{ NP[\str]$_j$ }\\[2mm]
%             \comps & \sliste{ NP[\type{lacc}]$_k$ }\\[2mm]
%             \argst & \sliste{ NP[\str]$_j$, NP[\type{lacc}]$_k$ }}}
%           [{\ms[stem]{
%             \phon  & \phonliste{ give }\\[2mm]
%             \da    & \sliste{ NP[\str]$_i$ }\\[2mm]
%             \argst & \sliste{ NP[\str]$_i$, NP[\str]$_j$, NP[\type{lacc}]$_k$ }}}]]
% \end{forest}
\hfill
%\scalebox{.8}{%
\begin{forest}
%sm edges
[{\ms[word]{\phon  & \phonliste{ gesägt }\\[2mm]
            \da    & \sliste{ NP[\str]$_i$ }\\[2mm]
            \argst & \sliste{ }}}
       [{\ms[stem]{
         \phon  & \phonliste{ säg }\\[2mm]
         \da    & \sliste{ NP[\str]$_i$ } \\[2mm]
         \argst & \sliste{ NP[\str]$_i$ } }}]]
\end{forest}%}
\hfill
%\scalebox{.8}{%
\begin{forest}
%sm edges
[{\ms[word]{\phon  & \phonliste{ zersägt }\\[2mm]
            \da    & \sliste{ NP[\str]$_i$ }\\[2mm]
%            comps  & \sliste{ NP[\str]$_j$ }\\
            \argst & \sliste{ NP[\str]$_j$ }}}
   [{\ms[stem]{\phon  & \phonliste{ zersäg }\\[2mm]
               \da    & \sliste{ NP[\str]$_i$ }\\[2mm]
               \argst & \sliste{ NP[\str]$_i$, NP[\str]$_j$ }}}
          [{\ms[stem]{
            \phon  & \phonliste{ säg }\\[2mm]
            \da    & \sliste{ NP[\str]$_i$ } \\[2mm]
            \argst & \sliste{ NP[\str]$_i$ } }}]]]
\end{forest}%}
\hfill\mbox{}
\caption{Anwendung der Passiv-Lexikonregel als Baum dargestellt, links für ein intransitives
  Verb, rechts im Zusammenspiel mit der \prefix{zer}-Präfigierungs-Lexikonregel, die ein
  intransitives auf ein transitives Verb abbildet}\label{fig-LR-as-tree}
\end{figure}

% 21.07.2024: seems to complicated.
Es gibt eine Konvention, dass Information, die in Lexikonregeln nicht erwähnt wird, von der Eingabe
unverändert in die Ausgabe der Lexikonregel übernommen wird. Für unser obiges Beispiel würde das
bedeuten, dass semantische Information unverändert übernommen wird. Siehe S.\,\pageref{le-gives} für die
Repräsentation der semantischen Beiträge.%
\label{page-lexical-rule-explanation-end}
% The bottom-most
% nodes in the tree (the leaves) are lexical entries for stems and the participle forms are derived by
% lexical rules and the new phonology values are provided at the licensed lexical items.

Wenden wir uns einigen prototypischen Verben und den Ergebnissen der Passivierungs-Lexikonregel zu: Die
\argstl des lizenzierten Partizips ist entweder leer (\mex{1}a) oder beginnt mit einem Objekt der
Aktivform:
\ea
\label{partizipien-hm}
%\resizebox{\linewidth}{!}{%
\begin{tabular}[t]{@{}l@{ }ll@{ }l@{}}
  &                              & \textsc{arg-st}\\[2mm]
a.&getanzt     (unerg):  & \liste{}\\[2mm]
b.&gelesen     (trans):    & \liste{ NP[\type{str}] }\\[2mm]
c.&geschenkt   (ditrans): & \liste{ NP[\type{ldat}], NP[\type{str}] }\\[2mm]
d.&geholfen    (unerg):  & \liste{ NP[\type{ldat}] }\\
\end{tabular}
%}
\z
Wie oben erläutert wurde, bekommt das erste Element in der \argstl mit strukturellem Kasus Nominativ, und
daher wird das Akkusativobjekt von \emph{lesen} in (\mex{1}a) in (\mex{1}b) als Nominativ realisiert:
\eal
\ex
Er liest den Aufsatz.
\ex
Der        Aufsatz wurde  gelesen.
\zl

\noindent
Das \ili{Englisch}e unterscheidet sich vom \ili{Deutsch}en darin, dass es überhaupt keinen Dativ hat. Ich spreche hier von morphologischen
Markierungen, nicht von Semantik. Daher sind beide Objekte ditransitiver Verben des \ili{Englisch}en
Akkusativobjekte. Allerdings kann nur eines der Objekte zum Subjekt befördert werden. Dies wird in
der vorliegenden Analyse dadurch modelliert, dass das sekundäre Objekt lexikalischen Akkusativ trägt (siehe auch \citew[\page 57]{Grewendorf2002a} für die Annahme eines lexikalischen Akkusativs für das sekundäre Objekt im \ili{Englisch}en).\footnote{
  Zugegebenermaßen ist dies nur eine Wiederholung der Fakten, da die Zuweisung von lexikalischem Kasus bedeutet, dass das
  betrachtete Argument keinen anderen Kasus haben kann. Aber zusammen mit Beschränkungen für
  Subjekte im \ili{Englisch}en folgen die Fakten über die Beförderung bzw.\ Nicht-Beförderung von Argumenten schön.
}

\ea\label{da-repr-hm-English}
%\resizebox{\linewidth}{!}{%
\begin{tabular}[t]{@{}l@{ }l@{ }l@{ }l@{ }l@{}}
  &                     & \textsc{arg-st}\\[2mm]
b.&dance   (unerg):     & \liste{ NP[\type{str}]}\\[2mm]
%c.&auf"|fallen (unacc): & \liste{}                         & \liste{NP[\type{str}], NP[\type{ldat}]}\\[2mm]
c.&read      (trans):   & \liste{ NP[\type{str}], NP[\type{str}]}\\[2mm]
d.&give      (ditrans): & \liste{ NP[\type{str}], NP[\type{str}], NP[\type{lacc}] }\\[2mm]
e.&help      (trans):   & \liste{ NP[\type{str}], NP[\type{str}] }\\
\end{tabular}
%}
\z
Das \ili{Deutsch}e kann das zweite Objekt (Akkusativ) befördern, das \ili{Englisch}e das erste Objekt. Die Gemeinsamkeit ist, dass
das verbnähere Objekt befördert werden kann. Im \ili{Deutsch}en ist das der Akkusativ, da
Nominativ"=Dativ"=Akkusativ die unmarkierte Abfolge ist und das \ili{Deutsch}e eine OV-Sprache ist, im \ili{Englisch}en aber der erste
Akkusativ, da das \ili{Englisch}e eine VO-Sprache ist.

\eal
\ex 
dass dem Kind der Ball gegeben wurde
\ex because the child was given the ball
\zl
Ein weiterer Unterschied ist der Lexikoneintrag für \emph{help}. Da es im \ili{Englisch}en keinen Dativ gibt, wird das
Objekt als Akkusativ markiert, wie es bei \emph{read} der Fall ist. Wie erwartet erlaubt das \ili{Englisch}e das
persönliche Passiv von \emph{help}, während dies im \ili{Deutsch}en nicht möglich ist:
\eal
\ex[]{
because he was helped
}
\ex[]{
weil ihm geholfen wurde
}
\ex[*]{
weil    er geholfen wurde
}
\zl

%% 
%% \subsubsubsection{Isländisch}

%% \begin{itemize}
%% \item Dativ und Genitiv sind lexikalisch:

%% \end{itemize}

%% }

\subsubsection{Primäre und sekundäre Objekte}

In diesem Abschnitt möchte ich Sprachen betrachten, die erlauben, dass beide Objekte befördert werden. Das \ili{Dänisch}e ist wie
das \ili{Englisch}e darin, dass es keinen Dativ hat. Das spiegelt sich in den folgenden \argstvs wider:
\ea\label{da-repr-hm-Danish}
%\resizebox{\linewidth}{!}{%
\begin{tabular}[t]{@{}l@{ }l@{ }l@{ }l@{ }l@{}}
  &                     & \textsc{arg-st}\\[2mm]
a.&danse   (tanzen, unerg):    & \liste{ NP[\type{str}]$_i$ }\\[2mm]
b.&læse      (lesen, trans):   & \liste{ NP[\type{str}]$_i$, NP[\type{str}]$_j$ }\\[2mm]
c.&give      (geben, ditrans): & \liste{ NP[\type{str}]$_i$, NP[\type{str}]$_j$, NP[\type{str}]$_k$ }\\[2mm]
d.&hjælpe    (helfen, trans):   & \liste{ NP[\type{str}]$_i$, NP[\type{str}]$_j$ }\\
\end{tabular}
%}
\z
Das \ili{Dänisch}e hat also zwei Objekte mit strukturellem Kasus, während das \ili{Englisch}e und das \ili{Deutsch}e nur ein Objekt mit strukturellem
Kasus haben und das andere mit lexikalischem Akkusativ bzw.\ lexikalischem Dativ. Da das \ili{Englisch}e und das
\ili{Deutsch}e keine Subjekte mit lexikalischem Kasus erlauben, ist klar, dass die Beförderung des Arguments mit lexikalischem Kasus zum Subjekt
ausgeschlossen ist. Das \ili{Dänisch}e erlaubt ebenfalls keine Subjekte mit lexikalischem Kasus, aber da die
zwei Objekte ohnehin strukturellen Kasus haben, können beide befördert werden.

Man beachte jedoch, dass die Lexikonregel in (\ref{lr-passive-prelim}) die Beförderung
des sekundären Objekts nicht erfasst. Was sie tut, ist, das Subjekt zu unterdrücken. Unter der Annahme, dass die
erste NP mit strukturellem Kasus das Subjekt ist, könnte das sekundäre Objekt niemals als
Subjekt realisiert werden. Man beachte, dass es nicht helfen würde zu sagen, dass jede NP mit strukturellem Kasus das Subjekt sein kann, da
dies falsche Realisierungen zuließe. Zusätzlich zum korrekten
(\ref{ex-fordi-manden-giver-barnet-bolden}) würden die folgenden beiden Sätze zugelassen:
\eal
\ex[*]{
\gll fordi barnet giver manden bolden\\ 
     weil Kind.\textsc{def} gibt Mann.\textsc{def} Ball.\textsc{def}\\\danish
}
\ex[*]{
\gll fordi   bolden         giver manden        barnet\\ 
     weil Ball.\textsc{def} gibt Mann.\textsc{def} Kind.\textsc{def}\\
}
\zl
(\mex{0}a) ist ungrammatisch mit \emph{barnet} `Kind' als Rezipient des Gebens. Ebenso kann das
transferierte Objekt \emph{bolden} in Aktivsätzen nicht als Subjekt realisiert werden. Das bedeutet, dass
die Beförderung zum Subjekt Teil der Lexikonregel sein muss, die das im Passiv verwendete Partizip lizenziert. Die Lexikonregel in (\mex{1}) nimmt die \argstl in der Eingabe der Lexikonregel und teilt sie in zwei
Listen auf: \ibox{1} und \ibox{2}. Die erste Liste \ibox{1} ist identisch mit dem Wert von \da. Die zweite Liste \ibox{2} ist der
Rest der \argstl. \ibox{2} ist mit \ibox{3}, dem \argstv der Ausgabe der Lexikonregel, durch die relationale Beschränkung
\texttt{promote} verbunden. \ibox{3} ist entweder gleich \ibox{2} oder eine Liste, in der eine weitere NP mit
strukturellem Kasus an den Anfang der Liste gesetzt wird.

\eas
\label{lr-passive-double-object}
Lexikonregel\is{Lexikonregel!Passiv} für das Passiv im \ili{Dänisch}en, \ili{Englisch}en, \ili{Deutsch}en und \ili{Isländisch}en:\\
\ms{
head   & \ms[verb]{ da & \ibox{1}\\
                  }\\
arg-st & \ibox{1} $\oplus$ \ibox{2} \\
} $\mapsto$
\ms{
arg-st & \ibox{3} \\
} $\wedge$ \texttt{promote}(\ibox{2}, \ibox{3})
\zs

%\largerpage
\noindent
(\mex{1}) zeigt die \argstvs unserer prototypischen Verben:
\ea\label{da-repr-hm-Danish-participles}
%\resizebox{\linewidth}{!}{%
\begin{tabular}[t]{@{}l@{ }l@{ }l}
  &                        & \textsc{arg-st}\\[2mm]
%a.&ankomme (unacc):       & \liste{}                         & \liste{NP[\type{str}]}\\[2mm]
a.&danset/-s   (tanzen, unerg):     & \liste{}\\[2mm]
%c.&auf"|fallen (unacc): & \liste{}                         & \liste{NP[\type{str}], NP[\type{ldat}]}\\[2mm]
b.&læst/-s      (lesen, trans):   &  \liste{NP[\type{str}]$_j$ } \\[2mm]
c.&givet/-s      (geben, ditrans): & \liste{NP[\type{str}]$_j$, NP[\type{str}]$_k$ } \\[2mm]
  &                         & \liste{NP[\type{str}]$_k$, NP[\type{str}]$_j$ } \\[2mm]
d.&hjulpet/-s    (helfen, trans):   & \liste{NP[\type{str}]$_j$ }                    \\
\end{tabular}
%}
\z
Die NP[\str]$_i$, die das erste Element in (\ref{da-repr-hm-Danish}) ist, wird unterdrückt. Der Effekt von
\texttt{promote} ist, dass es zwei verschiedene \argstvs für die Passivvarianten von \emph{givet}
`geben' gibt: eine mit einer \argstl, in der NP[\type{str}]$_j$ vor NP[\type{str}]$_k$ steht, und eine
andere, in der NP[\type{str}]$_j$ auf NP[\type{str}]$_k$ folgt. Die erste Abfolge entspricht
(\ref{ex-child-was-given-ball-danish}) -- hier wiederholt als (\ref{ex-child-was-given-ball-danish-two}) -- und
die zweite entspricht (\ref{ex-ball-was-given-child-danish}) -- hier wiederholt als (\ref{ex-ball-was-given-child-danish-two}):
\eal
\ex\label{ex-child-was-given-ball-danish-two}
\gll fordi barnet bliver givet bolden\\ 
     weil Kind.\textsc{def} wird gegeben Ball.\textsc{def}\\\danish
\glt `weil dem Kind der Ball gegeben wird'
\ex\label{ex-ball-was-given-child-danish-two}
\gll fordi bolden bliver givet barnet\\ 
     weil Ball.\textsc{def} wird gegeben Kind.\textsc{def}\\
\glt `weil der Ball dem Kind gegeben wird'
\zl
Bevor ich mich im nächsten Unterabschnitt den unpersönlichen Passiven im \ili{Dänisch}en zuwende, bespreche ich das Passiv in
Doppelobjektkonstruktionen im \ili{Isländisch}en.

Die Verteilung von strukturellem/lexikalischem Kasus im \ili{Isländisch}en ist im Grunde dieselbe wie im \ili{Deutsch}en. Der
Unterschied ist, dass das \ili{Isländisch}e Subjekte mit lexikalischem Kasus erlaubt und das \ili{Deutsch}e nicht.
(\mex{1}) zeigt unsere Standardbeispiele im \ili{Isländisch}en:
\ea\label{da-repr-hm-Icelandic}
%\resizebox{\linewidth}{!}{%
\begin{tabular}[t]{@{}l@{ }l@{ }l@{ }l@{ }l@{}}
  &                     & \textsc{arg-st}\\[2mm]
a.&dansa   (tanzen, unerg):     & \liste{ NP[\type{str}] }\\[2mm]
%c.&auf"|fallen (unacc): & \liste{}                         & \liste{NP[\type{str}], NP[\type{ldat}]}\\[2mm]
b.& lesa      (lesen, trans):   & \liste{ NP[\type{str}], NP[\type{str}] }\\[2mm]
c.&gefa       (geben, ditrans): & \liste{ NP[\type{str}], NP[\type{ldat}], NP[\type{str}] }\\[2mm]
d.&hjálpa     (helfen, trans):   & \liste{ NP[\type{str}], NP[\type{ldat}] }\\
\end{tabular}
%}
\z
Die Lexikonregel in (\ref{lr-passive-double-object}) lizenziert die folgenden Partizipien:

%\largerpage
\ea\label{da-repr-hm-Icelandic-two}%
%\scalebox{.98}{%
\resizebox{\linewidth}{!}{%
\begin{tabular}[t]{@{}l@{~}l@{~}l@{~~}l@{~~}l@{}}
  &                        & \textsc{arg-st}                     & \spr   & \comps\\[2mm]
%a.&ankomme (unacc):       & \liste{}                         & \liste{NP[\type{str}]}\\[2mm]
a.& dansað    (getanzt, unerg):     & \liste{}                        & \liste{ } & \liste{} \\[2mm]
%c.&auf"|fallen (unacc): & \liste{}                         & \liste{NP[\type{str}], NP[\type{ldat}]}\\[2mm]
b.& lesið      (gelesen, trans):   &  \liste{NP[\type{str}]$_j$ }         & \liste{NP[\type{str}]$_j$ } & \eliste\\[2mm]
c.& gefið      (gegeben, ditrans): & \liste{NP[\type{ldat}]$_j$, NP[\type{str}]$_k$ } & \liste{NP[\type{ldat}]$_j$ } & \liste{NP[\type{str}]$_k$ }\\[2mm]
  &                      & \liste{NP[\type{str}]$_k$, NP[\type{ldat}]$_j$ } & \liste{NP[\type{str}]$_k$ } & \liste{NP[\type{ldat}]$_j$ }\\[2mm]
d.& hjálpað    (geholfen, trans):   & \liste{NP[\type{ldat}]$_j$ }                  & \liste{ NP[\type{ldat}]$_j$ } & \liste{}\\
\end{tabular}%
}
\z
Zusätzlich zur \argstl zeigt (\mex{0}) die Abbildung auf die Merkmale \spr und \comps. Da
das \ili{Isländisch}e quirky Subjekte erlaubt, kann das Dativargument von `helfen' auf die
\sprl abgebildet werden \citep[\page 147--148]{Wechsler95a-u}. Ebenso ergeben die zwei Abfolgen der \argst von `geben' Partizipien mit einem dativischen
Subjekt und einem nominativischen Subjekt, wie es für die Analyse von (\ref{ex-were-the-king-given-the-slaves}) und (\ref{ex-were-the-slaves-given-the-king}) erforderlich ist, die hier
als (\mex{1}) wiederholt werden:
\eal
\ex\label{ex-were-the-king-given-the-slaves-two}
\gll Voru konunginum gefnar ambáttir?\\
     waren der.König.\DAT{} gegeben Sklavinnen.\NOM{}\\\icelandic
\glt `Wurden dem König Sklavinnen gegeben?'
\ex\label{ex-were-the-slaves-given-the-king-two}
\gll Var ambáttin gefin konunginum?\\
     war die.Sklavin.\NOM{} gegeben der.König.\DAT\\
\glt `Wurde die Sklavin dem König gegeben?'
\zl

\noindent
Das unpersönliche Passiv mit `tanzen' ist parallel zum unpersönlichen Passiv des \ili{Deutsch}en, aber die
Passivierung von `helfen' unterscheidet sich, da dies im \ili{Isländisch}en ein Fall des persönlichen Passivs ist.

\subsubsection{Unpersönliches Passiv}
\label{sec-impersonals}

%\largerpage
Als letzten Punkt in diesem Unterabschnitt werfen wir einen Blick auf das unpersönliche Passiv. Das \ili{Deutsch}e und das
\ili{Isländisch}e verlangen keine Subjekte. Wenn es also keine NP mit strukturellem Kasus gibt, ist die Konstruktion im
\ili{Deutsch}en subjektlos. Ebenso verlangt das \ili{Isländisch}e kein Subjekt: Wenn es kein NP-Argument gibt,
ist das Ergebnis ein unpersönliches Passiv. Ein Beispiel für den letzteren Fall ist die Passivierung von
\emph{dansa} `tanzen'. Die \argstl ist die leere Liste, und daher sind die \sprl und die \compsl
ebenfalls leer. Passivpartizipien von Verben, die ein NP- und ein PP-Objekt regieren, haben eine \argstl,
die nur das PP-Argument enthält. Dieses PP-Argument wird auf die \compsl abgebildet, und somit ergibt sich eine
subjektlose Konstruktion.

Das \ili{Englisch}e erlaubt keine unpersönlichen Passive, da es eine NP oder ein satzwertiges Argument verlangt, das
als Subjekt dienen kann. Das \ili{Dänisch}e verlangt ebenfalls ein Subjekt, erlaubt aber unpersönliche
Konstruktionen. Der Trick, den das \ili{Dänisch}e anwendet, ist die Einfügung eines
Expletivums\is{Pronomen!expletives}. Das Interessante ist, dass das Expletivum bei allen
Verben eingefügt werden kann, die kein Objekt regieren. Das heißt, Passive wie das Beispiel in (\ref{ex-at-der-bliver-danset}) und intransitive
Verben wie die in (\mex{1}) können ein expletives Subjekt haben, und alle anderen Argumente werden
postverbal realisiert.
\eal
% Bjarne 26.07.2023
\ex 
\gll at    der    ikke går   en mand på gaden\\
     dass \expl{} nicht geht ein Mann auf Straße.\textsc{def}\\\danish
\ex 
\gll at   der     ikke kommer to  nye medarbejdere\\
     dass \expl{} nicht kommen zwei neue Mitarbeiter\\
\ex
\gll at   der     ikke venter nogle hårde forhandlinger\\
     dass \expl{} nicht warten einige harte Verhandlungen\\
\zl
Die Beispiele in (\mex{0}) sind eingebettete Sätze mit Negationen, was zeigt, dass das Expletivum nicht
in einem V2-Satz vorangestellt ist, sondern dass wir hier tatsächlich ein Expletivum in Subjektposition haben.

% Das sollte eine Lexikonregel sein, die das Expletivum einführt. Die Idee wäre, dass man den
% Lexikoneintrag ohne Expletivum hat, dass man den aber nicht benutzen kann. Das Problem ist aber,
% dass man dann die Forderung nach dem Subjekt nicht in das Mapping einbauen kann, denn dann wäre
% der Zwischenlexikoneintrag einfach schon inkonsistent und könnte nicht als Input für die
% Expletivregel genommen werden.

In einer früheren Auflage dieses Buches bin ich davon ausgegangen, dass die
Expletiveinfügung während der Abbildung der \argst-Elemente auf \spr und \comps stattfindet. Wenn es
eine NP/VP/CP am Anfang der \argstl gibt, wird sie auf \spr abgebildet, und alle anderen Elemente werden
auf \comps abgebildet. Wenn es kein Objekt gibt, wird ein Expletivum eingefügt. Das erfasst sowohl
subjektlose Passive als auch Expletivkonstruktionen von Intransitiva wie (\mex{1}).\footnote{%
  \citet[Abschnitt~4.3]{BB2007a} nehmen an, dass dänische Witterungsverben\is{Verb!Witterungs-} keine Argumente auf ihrer
  \argstl haben, sondern dass ein Expletivum in der Abbildung von \argst auf die Valenzmerkmale eingefügt wird. Ohne
  Details anzugeben, schlagen sie vor, dass sich dies auch auf die oben besprochenen Fälle erstreckt.
}

%\largerpage
%\enlargethispage{5pt}
(\mex{1}) zeigt die Abbildungen, die sich in diesem Ansatz für das \ili{Dänisch}e ergaben.
\ea\label{da-repr-hm-Danish-three}
%\resizebox{\linewidth}{!}{%
\begin{tabular}[t]{@{}l@{ }l@{ }l@{ }l@{ }l@{~~~~~}l@{}}
  &                        & \textsc{arg-st}                     & \spr   & \comps\\[2mm]
%a.&ankomme (unacc):       & \liste{}                         & \liste{NP[\type{str}]}\\[2mm]
%a.&danset/-s   (unerg):     & \liste{}                        & \liste{ NP$_\typeSM{expl}$ } &
%\liste{} \\[2mm]
a.&danset/-s   (unerg):     & \liste{}                        & \liste{ NP\textsubscript{\normalfont\itshape expl\/} } & \liste{} \\[2mm]
%c.&auf"|fallen (unacc): & \liste{}                         & \liste{NP[\type{str}], NP[\type{ldat}]}\\[2mm]
b.&læst/-s      (trans):   &  \liste{ NP[\type{str}]$_j$ }                     & \liste{ NP[\type{str}]$_j$ } & \eliste\\[2mm]
c.&givet/-s      (ditrans): & \liste{ NP[\type{str}]$_j$, NP[\type{str}]$_k$ } & \liste{ NP[\type{str}]$_j$ } & \liste{ NP[\type{str}]$_k$ }\\[2mm]
  &                         & \liste{ NP[\type{str}]$_k$, NP[\type{str}]$_j$ } & \liste{ NP[\type{str}]$_k$ } & \liste{ NP[\type{str}]$_j$ }\\[2mm]
d.&hjulpet/-s    (trans):   & \liste{ NP[\type{str}]$_j$ }                     & \liste{ NP[\type{str}]$_j$ } & \liste{ }\\
\end{tabular}
%}
\z

\noindent
Das Problem bei einem solchen Ansatz ist, dass die Abbildung von \argst auf die Valenzmerkmale
einzelsprachspezifisch und ziemlich kompliziert wäre. Stattdessen scheint es angemessener, Lexikonregeln
für die Einfügung von Expletiva anzunehmen. Ich nehme an, dass Verben, die kein direktes Objekt oder gar kein
Argument haben, Eingabe der Lexikonregel sein können. Die Lexikonregel fügt das Expletivum hinzu, und dieses kann
dann auf die \sprl abgebildet werden \citep{MuellerDanishImpersonal}. (\mex{1}) zeigt die Abbildung der
Version von \emph{danset} `getanzt', die diese Lexikonregel durchlaufen hat:
\ea
\begin{tabular}[t]{@{}l@{ }l@{ }l@{ }l@{ }l@{~~~~~}l@{}}
  &                        & \textsc{arg-st}                     & \spr   & \comps\\[2mm]
  &danset/-s   (unerg):     & \liste{ NP\textsubscript{\normalfont\itshape expl\/} }                        & \liste{ NP\textsubscript{\normalfont\itshape expl\/} } & \liste{} \\[2mm]
\end{tabular}
\z

\subsubsection{Abbildung in den germanischen OV-Sprachen}
\label{sec-Mapping in the Germanic OV languages}

Die germanischen OV-Sprachen unterscheiden sich von den VO-Sprachen darin, dass sie mit dem
Hilfsverb einen Verbalkomplex bilden. Das hat Konsequenzen für den Rest des grammatischen Systems, etwa für
Voranstellungen. (\mex{1}) zeigt die Abbildungen von \argst auf \spr und \comps für prototypische deutsche Verben:
\ea
\label{partizipien-hm}%
\resizebox{.99\linewidth}{!}{%
\begin{tabular}[t]{@{}l@{ }l@{~~}l@{ }l@{~~}l@{}}
  &                              & \textsc{arg-st} & \spr & \comps\\[2mm]
a.&getanzt     (unerg):  & \liste{}        & \liste{} & \liste{}\\[2mm]
b.&gelesen     (trans):    & \liste{ NP[\type{str}] } & \liste{} & \liste{ NP[\type{str}] }\\[2mm]
c.&geschenkt   (ditrans): & \liste{ NP[\type{ldat}], NP[\type{str}] } & \liste{} & \liste{ NP[\type{ldat}], NP[\type{str}] }\\[2mm]
d.&geholfen    (unerg):  & \liste{ NP[\type{ldat}] } & \liste{} & \liste{ NP[\type{ldat}] }\\
\end{tabular}}
\z
Man beachte, dass beide Mitglieder der \argstl auf \comps abgebildet werden. Der \compsv ist also eher derjenige, den man
beim Perfektpartizip vermuten würde. Dies wird in Abschnitt~\ref{sec-perfect} ausführlicher diskutiert.
Für den Moment möchte ich darauf hinweisen, dass jede Teilmenge der Komplemente nicht-finiter Verben zusammen
mit dem Verb vorangestellt werden kann:
\eal
\ex 
Den        Wählern Märchen erzählen sollte man besser nicht.
\ex 
Den        Wählern erzählen sollte man solche Märchen     nicht.
\ex 
Märchen     erzählen sollte man den        Wählern nicht.
\zl
Das Subjekt ist nicht Teil der \compsl von Infinitiven mit oder ohne \emph{zu}, und daher sagt die
Analyse in Kapitel~\ref{chap-verbal-complex} voraus, dass Beispiele wie (\mex{1})
ungrammatisch sind:
\eal
\ex[*]{
Dieser Mann erzählen sollte den Wählern solche Geschichten besser nicht.
}
\ex[*]{ 
Dieser      Mann den Wählern       erzählen sollte solche Geschichten besser nicht.
}
\ex[*]{
Dieser      Mann solche Geschichten erzählen sollte den        Wählern besser nicht.
}
\zl
Bei den Partizipien ist die Situation anders: Unabhängig von der Frage, ob wir eine
Perfekt- oder eine Passivkonstruktion haben, kann die NP, die im Aktiv das Akkusativobjekt und im
Passiv das Subjekt ist, vorangestellt werden:
\eal
\ex 
Diesen      Plan den        Wählern gegeben hat er nicht.
\ex 
Dieser Plan den Wählern gegeben wurde damals nicht.
\zl
Dies wird erfasst, indem alle nicht unterdrückten Argumente von Partizipien auf \comps abgebildet werden. Aus Platzgründen wird das Phänomen der partiellen Voranstellung von Verbalphrasen hier nicht behandelt. Der interessierte
Leser sei auf \citew{Mueller96a}, \citew[Kapitel~2.2.2]{Mueller2002b} bzw.\ \citew[Kapitel~15.2]{MuellerLehrbuch4} verwiesen.

%\if 0
%\subsection{Variation and Generalizations}

\subsection{Das Passivhilfsverb}
\label{sec-auxiliary}

%\largerpage
Wir haben nun gesehen, wie die Partizipformen der betrachteten Sprachen aussehen und wie sie
aus Lexikoneinträgen für Stämme über Lexikonregeln lizenziert werden. Was fehlt, sind die Lexikoneinträge für die
Hilfsverben. Kapitel~\ref{chap-verbal-complex} hat sich mit der Analyse von Verbalkomplexen in SOV-Sprachen wie
dem \ili{Deutsch}en, \ili{Niederländisch}en und \ili{Afrikaans} befasst, und es wurde darauf hingewiesen, dass SVO-Sprachen wie das \ili{Englisch}e und die
\ili{skandinavisch}en Sprachen keinen Verbalkomplex bilden. Vor diesem Hintergrund mag es überraschen,
dass es möglich ist, eine allgemeine Beschränkung zu formulieren, die alle Passivhilfsverben in den
\ili{germanisch}en Sprachen erfasst. Das folgende AVM zeigt eine Beschränkung, die für alle \argst-Listen für Passivhilfsverben in \ili{germanisch}en Sprachen gilt:\footnote{%
Der von \citet[\page 147]{Mueller2002b} und \citet[\page 149]{MOe2013a} verwendete Lexikoneintrag des Passivhilfsverbs
spezifiziert den \dav des eingebetteten Partizips als referentielle NP. Das schließt die
Passivierung unakkusativischer\is{Verb!unakkusativisches} Verben aus, die die leere Liste als \dav haben.
}
 
\ea
\label{le-passive-aux-arg-st}
Passivhilfsverb für die \ili{germanisch}en Sprachen (\emph{be}, \emph{werden}, \emph{blive} etc.):\\
\avm{
[ arg-st \1 \+ \2 \+  < [ vform & ppp\\
%                          da    & < XP$_{ref}$ >\\
                          spr   & \1\\
                          comps & \2  ] > ] 
}
\z
%
% The specification of the \dav excludes unaccusative verbs from forming a passive since unaccusative verbs have the empty
% list as the value of their \textsc{da} feature (the \textsc{da} feature has as its value a list
% containing the argument with subject properties). Weather verbs have an element in their \textsc{da}
% list but it is non-referential. Specifying the \textsc{da} element to be referential excludes the
% passive of weather verbs like \emph{regnen} `to rain', which otherwise could enter an impersonal
% passive:
% \ea[*]{
% \gll weil dann doch geregnet wurde\\
%      because then nevertheless rained was\\
% \glt Intended: `because there was raining nevertheless'
% }
% \z

%\largerpage
%\enlargethispage{5pt}
\noindent
Wenn das Passivhilfsverb in (\ref{le-passive-aux-arg-st}) in einer Grammatik des \ili{Deutsch}en verwendet wird, werden die
Argumente des Partizips (\ibox{1} und \ibox{2}) vom Passivhilfsverb angezogen
\citep{HN89a,HN94a}.
%The Verbal Complex Schema allows the combination of the auxiliary with an unsaturated
%non-head daughter. The SVO languages do not have the verbal complex schema. They 
Hilfsverben in den SVO-Sprachen betten eine VP ein. Das bedeutet, dass der \compsv der eingebetteten verbalen
Projektion die leere Liste sein muss. Daher ist \ibox{2} die leere Liste, und nur der Spezifikator \iboxb{1} wird
vom eingebetteten Verb übernommen.

%% \item Hence, we have explained how
%% \ili{Danish} and \ili{English} embed a VP and \ili{German} forms a verbal complex although the lexical item of the
%% auxiliary does not require a VP complement.
%\fi

%\largerpage
Mit dem nun vorliegenden Lexikoneintrag für das Hilfsverb können wir uns einige Beispielanalysen von
Passivsätzen ansehen. Beginnen wir mit einem \ili{englisch}en Beispiel. Die Analyse von (\mex{1}) ist in
Abbildung~\ref{fig-the-child-was-given-a-novel} dargestellt.
\ea
The child was given a novel.
\z
\begin{figure}
\scalebox{1}{%
\begin{forest}
sm edges
[V\feattab{
            \spr    \sliste{ },\\
            \comps  \sliste{ }}
   [\ibox{1} \npnom [the child,roof]]
   [V\feattab{
            \spr    \sliste{ \ibox{1} },\\
            \comps  \sliste{ }}
      [V\feattab{
%            \da     \sliste{ \ibox{2} },\\
            \spr    \sliste{ \ibox{1} },\\
            \comps  \sliste{ \ibox{2} },\\
            \argst \sliste{ \ibox{1}, \ibox{2} }}  [was]]
      [\ibox{2} V\feattab{
%           \da     \sliste{ \ibox{3} },\\
            \spr    \sliste{ \ibox{1} },\\
            \comps  \sliste{ }}, s sep+=1em
        [V\feattab{
%            \da     \sliste{ \ibox{3} },\\
            \phon   \phonliste{ given },\\
            \spr    \sliste{ \ibox{1} },\\
            \comps  \sliste{ \ibox{3} },\\
            \argst \sliste{ \ibox{1}, \ibox{3} }} 
          [V\feattab{
            \da    \ibox{4} \sliste{ NP[\str] },\\
            \argst \ibox{4} $\oplus$ \sliste{ \ibox{1}, \ibox{3} }} [{give-}]]]
        [\ibox{3} \npacc [a novel,roof]]]]]
\end{forest}}
\caption{\label{fig-the-child-was-given-a-novel}Analyse von \emph{The child was given a novel.}}
\end{figure}

\noindent
Der Lexikoneintrag für den Stamm \stem{give} ist Eingabe der Passiv-Lexikonregel. Die Passiv-Lexikonregel
lizenziert das Partizip \emph{given}. Die \argst von \stem{give} wird um das Element in
der \da-Liste von \stem{give} \iboxb{4} verkürzt. Das Ergebnis ist eine \argstl, die die beiden NPs enthält, die in
Aktivsätzen die Objekte wären. Die erste \iboxb{1} wird auf \spr abgebildet und die zweite \iboxb{3}
auf \comps. Die Kombination von \emph{given} und \emph{a novel} bildet eine VP (etwas mit einer leeren
\compsl und einem Element in der \sprl). Das Passivhilfsverb \emph{was} selegiert die VP \emph{given a
  novel}.
%It requires the selected VP to have a referential NP as designated argument, which the VP
%has \iboxb{4}. 
Der Spezifikator der selegierten VP \iboxb{1} wird angezogen. Als erstes Argument von
\emph{was} mit strukturellem Kasus bekommt diese NP Nominativ. Schließlich wird die VP \emph{was given a novel}
mit \emph{the child} kombiniert, und wir haben einen vollständigen Satz.

Die Analyse des parallelen \ili{deutsch}en Satzes in (\mex{1}) ist in
Abbildung~\vref{fig-dem-Kind-ein-Roman-gegeben-wurde} dargestellt.
%\largerpage
\ea
dass dem Kind ein Roman gegeben wurde
\z
% Without this, the footnote floats to the next page. Hm.
%\pagebreak

\begin{figure}
\scalebox{.95}{%
\begin{forest}
sm edges
[V\feattab{
            \spr    \sliste{ },\\
            \comps  \sliste{ }}
   [\ibox{1} \npdat [dem Kind,roof]]
   [V\feattab{
            \spr    \sliste{  },\\
            \comps  \sliste{ \ibox{1} }}, s sep+=2em
      [\ibox{2} \npnom [ein Roman,roof]]
      [V\feattab{
%            \da     \sliste{ \ibox{3} },\\
            \spr    \sliste{ },\\
            \comps  \sliste{ \ibox{1}, \ibox{2} }}
          [\ibox{3} V\feattab{                                  % gegeben
%              \da    \sliste{ \ibox{3} },\\
              \phon  \phonliste{ gegeben },\\
              \spr   \sliste{ },\\
              \comps \sliste{ \ibox{1}, \ibox{2} },\\
              \argst \sliste{ \ibox{1}, \ibox{2} }} 
            [V\feattab{
              \da    \ibox{4} \sliste{ NP[\str]  },\\
              \argst \ibox{4} $\oplus$ \sliste{ \ibox{1}, \ibox{2} }} [{geb-}]]]
          [V\feattab{
%            \da    \sliste{ \ibox{2} },\\
             \spr   \sliste{  },\\
             \comps \sliste{ \ibox{1}, \ibox{2}, \ibox{3} },\\
             \argst \sliste{ \ibox{1}, \ibox{2}, \ibox{3} }}  [wurde]]]]]
\end{forest}}
\caption{\label{fig-dem-Kind-ein-Roman-gegeben-wurde}Analyse von \emph{dass dem Kind ein Roman gegeben
    wurde}}
\end{figure}

\noindent
Die Analyse ähnelt der des \ili{englisch}en Beispiels, unterscheidet sich aber darin, dass das Hilfsverb und
das Partizip einen Verbalkomplex bilden. Zunächst wird die Passiv-Lexikonregel auf einen Verbstamm
\stem{geb} angewendet und lizenziert die Partizipform \emph{gegeben}. \emph{gegeben} \iboxb{3} wird mit
\emph{wurde} kombiniert, und \emph{wurde} übernimmt die Elemente der \compsl von \emph{gegeben} \sliste{
  \ibox{1}, \ibox{2} }. Das Ergebnis der Kombination von \emph{gegeben} und \emph{werden} hat die
\compsl \sliste{ \ibox{1}, \ibox{2} }.

\subsection{Das morphologische Passiv}

%\largerpage
Eine Lexikonregel, die der für die Partizipformen ähnelt, kann für die
morphologischen Passive im \ili{Dänisch}en verwendet werden. Ein Unterschied ist natürlich das affixale Material, das von
der Regel hinzugefügt wird. Außerdem wird für das morphologische Passiv angenommen, dass das \da der Eingabe der Lexikonregel ein
referentielles XP enthalten muss. Wie im vorigen Abschnitt diskutiert wurde, schließt das morphologische Passive von
Unakkusativa\is{Verb!unakkusativisches} und Witterungsverben\is{Verb!Witterungs-} aus. 

%% \subsection{Agent Expressions}

%% We follow \citet[Chapter~7]{Hoehle78a} and \citet[Section~5]{Mueller2003e} and treat the \emph{by}
%% phrases as adjuncts.

\subsection{Das Perfekthilfsverb}
\label{sec-perfect}

\citegen{Haider86} Analyse\is{Perfekt|(} des Passivs ist brillant, da es ausreicht, einen Lexikoneintrag
für das Partizip zu haben. Das Partizip hat ein blockiertes designiertes Argument, und das designierte
Argument bleibt im Passiv blockiert, während das Perfekthilfsverb das designierte
Argument deblockiert. 
\eal
\ex
dass der        Aufsatz gelesen wurde
\ex
dass Kirby den Aufsatz gelesen hat
\zl

%\largerpage
\noindent
Die Deblockierung des designierten Arguments ist möglich, da das designierte Argument nicht nur im Stamm eines Verbs
kodiert ist, sondern auch im Lexikoneintrag für das Partizip.   

\ea
Lexikoneintrag für Perfekthilfsverben in SOV-Sprachen wie dem \ili{Niederländisch}en und \ili{Deutsch}en:\\
\avm{
[ arg-st \1 \+ \2 \+ \3 \+  < [ vform & ppp\\
                                da    & \1\\
                                spr   & \2\\
                                comps & \3 ] > ]
}
\z

%\subsubsection{Analyse als komplexes Prädikat für Dänisch und Englisch?}

\noindent
Leider funktioniert \citeauthor{Haider86}s Ansatz für SVO-Sprachen nicht. Wenn wir den
Ansatz der Argumentblockierung/-deblockierung verwenden wollten, müssten wir die Strukturen in (\mex{1}a--b) annehmen:
\eal
\ex He [has given] the book to Mary.
\ex The book [was given] to Mary.
\ex He has [given the book to Mary].
\ex The book was [given to Mary].
\zl
Wenn wir annehmen, dass Hilfsverben VPs einbetten, wie auf Seite~\pageref{page-English-Aux-VPs} argumentiert wurde, geraten wir
in Probleme, da das Subjekt des Partizips blockiert ist und die einzige VP, die wir mit dem
Partizip bilden können, die in (\mex{0}d) ist, aber für das Perfekt brauchen wir auch eine VP, die das Objekt enthält, wie
in (\mex{0}c).%
\is{Perfekt|)}

%% 
%% \subsubsection{A Solution that Almost Works}

%% \begin{itemize}
%% \item Complex Passive: There has to be a way to distinguish between participles that can be used in both perfect and
%%   passive:\\
%% \textsc{voice} feature. 

%% \begin{itemize}
%% \item Value is \type{passive} for those participles that cannot be used in perfect constructions.

%% 
%% \item Value is underspecified for participles that can be used in both perfect and passive

%% 
%% \item Perfect requires \textsc{voice} value to be \type{active}.
%% \end{itemize}

%% 
%% \item Expletives: Perfect attracts args from \argstl rather than \spr/\comps.
%% \begin{itemize}
%% \item Since expletives are not on \argst, they will not get into the way.
%% \end{itemize}
%% \end{itemize}

%% }

%\subsubsection{But: (Partial) Fronting}
%\subsubsection{Problem: (Partial) Fronting}

% \citet{Meurers99b} hat einen Trick gefunden, wie man die Kasuszuweisung in (\mex{1})
%   analysieren kann:
% \nocite{Meurers2000b,MdK2001a}
% \eal
% \ex 
% Gelesen wurde der Aufsatz schon oft.
% \ex 
% Der Aufsatz gelesen wurde schon oft.
% \ex
% Den Aufsatz gelesen hat er schon oft.
% \zl

% Das funktioniert aber nicht für Dänisch/Englisch, denn hier haben wir nicht nur Kasus- sondern
%   auch Positionsunterschiede:
% \eal
% \ex The book should have been given to Mary and\\
%     {}[given to Mary] it was.
% \ex He wanted to give the book to Mary and\\
%     {}[given the book to Mary] he has.
% \zl

% Wenn sich keine ausgeklügelten Mechanismen für die Unterspezifikation verschiedenenr Mappings finden
% lassen, müssen wir wohl zwei verschiedene Partizipformen annehmen.
% for the participle form.

\subsection{Das Fernpassiv}
\label{sec-remote-passive-phen}

%\largerpage
Das sogenannte \emph{Fernpassiv}\is{Passiv!Fernpassiv} ist ein Höhepunkt der Syntax des \ili{Deutsch}en, da mehrere
Phänomene auf nicht-triviale Weise interagieren. Es wurde zuerst von \citet[\page
175--176]{Hoehle78a} diskutiert. Höhle hat beobachtet, dass Objekte deutscher Infinitive mit \emph{zu} in bestimmten Kontexten im
Nominativ erscheinen. (\mex{1}) liefert einige konstruierte Beispiele aus der Literatur:
\eal
\ex\iw{versuchen|(}
daß  er        auch von mir zu überreden versucht wurde\footnote{
        \citew*[\page 212]{Oppenrieder91a}.%
}

\ex 
weil    der        Wagen oft   zu reparieren versucht wurde
\label{bsp-zu-reparieren-versucht-wurde}
\zl
Die Beispiele in (\mex{1}) sind belegte Daten, die von \citet[\page 136--137]{Mueller2002b} aus dem COSMAS-Korpus gesammelt wurden:
\eal
\ex Dabei darf jedoch nicht vergessen werden, daß in der Bundesrepublik, wo \emph{ein Mittelweg zu gehen versucht wird}, 
die Situation der Neuen Musik allgemein und die Stellung der Komponistinnen im besonderen noch recht unbefriedigend ist.\footnote{
Mannheimer Morgen, 26.09.1989, Feuilleton; Ist's gut, so unter sich zu bleiben?
}
%But should however not forgotten get that in the BRD where a middle.way to go tried gets the situation of.the 
%new music generally and the position of.the composers in particular still quite unsatisfactory is 
%
\glt `Man sollte nicht vergessen, dass die Situation der Neuen Musik im Allgemeinen und die Stellung der Komponistinnen im Besonderen in der Bundesrepublik, wo man einen Mittelweg zu gehen versucht, recht unbefriedigend ist.'
\ex Noch ist es nicht so lange her, da ertönten gerade aus dem Thurgau jeweils die lautesten Töne, 
    wenn im Wallis oder am Genfersee im Umfeld einer Schuldenpolitik mit den unglaublichsten Tricks 
    \emph{der sportliche Abstieg zu verhindern versucht wurde}.\footnote{
St.\ Galler Tagblatt, 09.02.1999, Ressort: TB-RSP; HCT und das Prinzip Hoffnung.%
}
%still is it not so long ago there sounded just out of.the Thurgau at.the.time the 
%loudest sounds when in.the Valais or at.the Lake.Geneva in.the sphere of.a debt.policy 
%with the most.unbelievable tricks the sporty relegation to prevent tried got
%
\glt `Es ist noch nicht allzu lange her, dass gerade aus dem Thurgau die lautesten Töne ertönten, als im Wallis oder am Genfersee im Umfeld einer Schuldenpolitik mit den unglaublichsten Tricks versucht wurde, den sportlichen Abstieg zu verhindern.' 
\ex Die Auf- und Absteigenden erzeugen ungewollt einen Ton, \emph{der bewusst nicht als lästig zu eliminieren versucht wird}, 
    sondern zum Eigenklang des Hauses gehören soll, so wünschen es sich die Architekten.\footnote{
Züricher Tagesanzeiger, 01.11.1997, p.\,61.%
}

% Philippa
\glt `Die hinauf- und hinabsteigenden Menschen erzeugen ungewollt einen Ton, den man bewusst nicht zu eliminieren versucht, der aber vielmehr zum Eigenklang des Hauses gehören soll, so wünschen es sich die Architekten.'
% Uta
%the up and downclimbers create involuntarily a tone that consciously not as annoying 
%to eliminate tried gets but to.the own.sound of.the house belong should so wish it themselves the architects
%`That no attempt is made to eliminate the involuntary noise caused by people ascending and descending 
%is a conscious decision; this is not considered to be a nuisance, but as an aspect of the house's own sound, 
%at least according to the architects.'
\zl
Höhles Beispiele und andere Beispiele aus der Literatur betrafen das Verb \emph{versuchen},
aber \citet{Wurmbrand2003a} hat gezeigt, dass auch andere Verben das Fernpassiv erlauben. (\mex{1})
und (\mex{2}) zeigen einige ihrer Beispiele mit \emph{beginnen}, \emph{vergessen} und
\emph{wagen}:
\ea
%% \ex
%% \emph{dieser} wurde bereits zu bauen begonnen.\footnote{
%%         \url{http://www.hollabrunn.noe.gv.at/mariathal/ortsvorsteher.html}, 28.07.2003.
%% }
\emph{der} \emph{zweite} \emph{Entwurf} wurde  zu bauen begonnen,\footnote{
\url{http://www.waclawek.com/projekte/john/johnlang.html}, 28.07.2003.
}

\z

\noindent
Während der Kasus von \emph{der zweite Entwurf} eindeutig Nominativ ist, ist dies bei
den Beispielen in (\mex{1}) nicht der Fall, da die jeweiligen Elemente im
Plural stehen und daher Nominativ oder Akkusativ sein könnten. Aber aufgrund der \isi{Kongruenz} mit dem finiten Verb ist
klar, dass die Relativpronomen im Nominativ stehen.
\eal
\ex 
Anordnungen, die zu stornieren vergessen \emph{wurden}\footnote{
        \url{http://www.rlp-irma.de/Dateien/Jahresabschluss2002.pdf}, 28.07.2003.
}

\ex Aufträge [\ldots], die zu drucken vergessen worden \emph{sind}\/\footnotemark\\
    Aufträge {} die zu drucken vergessen worden sind\\
\footnotetext{
        \url{http://www.iitslips.de/news.html}, 28.07.2003.
}
\glt `Aufträge, die zu drucken vergessen worden sind'
%\ex Ist plötzlich übervoll von Emotionen und längst begrabenen Träumen, die nicht zu leben gewagt wurden\footnote{
% nicht auffindbar
\ex NUR Leere, oder doch noch Hoffnung, weil aus Nichts wieder Gefühle entstehen,\\
die so vorher nicht mal zu träumen gewagt \emph{wurden}?\footnote{
        \url{http://www.ultimaquest.de/weisheiten_kapitel1.htm}, 28.07.2003.
}

\ex Dem Voodoozauber einer Verwünschung oder die gefaßte Entscheidung zu einer Trennung,\\
die bis dato noch nicht auszusprechen gewagt \emph{wurden}\footnote{
        \url{http://www.wedding-no9.de/adventskalender/advent23_shawn_colvin.html}, 28.07.2003.
}

\zl
% Kasus bei PVP wie Haiders entziffern: Am leichtesten zu erklären fiel den 
% Experten dabei gestern der Kursverlust der Telekom, zu deren Schuldenproblem 
% eine neue Hiobsbotschaft kam.  (taz. 8./9. 9. 01 S. 9.)
%

\noindent
Das Objekt eines Verbs, das unter einem Passivpartizip eingebettet ist, wird zum Subjekt des Satzes befördert:
\eal
\ex 
\longexampleandlanguage{
weil    Aicke den        Wagen oft   zu reparieren versucht hat
\ex 
weil    der        Wagen oft   \emph{zu} \emph{reparieren} \emph{versucht} \emph{wurde}
\label{bsp-zu-reparieren-versucht-wurde-two}
\zl

\noindent
Das Fernpassiv ist nur in Verbalkomplexen möglich. Wenn kein Verbalkomplex gebildet wird wie in
(\mex{1}a,c), muss das Objekt von \emph{reparieren} im Akkusativ erscheinen:
\eal
\ex[]{
weil    oft   versucht wurde, den        Wagen zu reparieren.
}
\ex[*]{
weil    oft   versucht wurde, der        Wagen zu reparieren.
}
\ex[]{
Den        Wagen zu reparieren wurde  oft   versucht.
}
\ex[*]{
Der        Wagen zu reparieren wurde  oft   versucht.
}
\zl
Der Unterschied zwischen (\ref{bsp-zu-reparieren-versucht-wurde-two}) und (\mex{0}a,c) wird durch eine Analyse erklärt, die das Fernpassiv als Passivierung
eines Prädikatskomplexes behandelt, \dash durch eine Analyse, die \pref{bsp-zu-reparieren-versucht-wurde-two} die Struktur (\mex{1}) zuweist.
\ea
weil    der        Wagen oft   [[zu reparieren versucht] wurde].
\label{bsp-zu-reparieren-versucht-wurde-three}
%
\z
In (\mex{-1}a,c) haben wir keine Prädikatskomplexe. Das Objekt von \emph{zu reparieren} ist Teil der VP
und bekommt daher Akkusativ. Die Passive in (\mex{-1}a,c) sind unpersönliche Passive\is{Passiv!unpersönliches}.

Das Verb \emph{versuchen} selegiert ein Subjekt, einen Infinitiv mit \emph{zu} und die
Komplemente des eingebetteten Verbs.
\ea
\stem{versuch}:\\
\avm{
[ arg-st & < !NP[\type{str}]$_i$! > \+ \1 \+ < V[\type{inf}, \subj < !NP[\type{str}]$_i$! >, \comps \1 ] > ]
}
\z

In unserem Beispiel ist das eingebettete Verb \emph{reparieren} und hat ein Komplement. (\mex{1}) zeigt die
\argstv von \stem{versuch}. Die erste NP ist das Subjekt von \emph{versuchen} und die zweite NP ist das
angezogene Objekt von \emph{zu reparieren}.
\ea
\argstv von \stem{versuch} bei Einbettung eines streng transitiven Verbs:\\
\sliste{ NP[\type{str}]$_i$, NP[\type{str}]$_j$, V[\type{inf}] }
\z
%\largerpage
Wenn die Passiv-Lexikonregel auf diesen Verbstamm angewendet wird und das jeweilige Partizip lizenziert, hat der
resultierende Lexikoneintrag für \emph{versucht} die folgende \argstv:
\ea
\argst des Partizips \emph{versucht} bei Einbettung eines streng transitiven Verbs:\\
\sliste{ NP[\type{str}]$_j$, V[\type{inf}] } 
\z

%\largerpage
\noindent
Die zwei verbleibenden \argst-Mitglieder werden auf \comps abgebildet (siehe Abschnitt~\ref{sec-Mapping in the Germanic OV languages}). Nach der Kombination mit
\emph{zu reparieren} haben wir einen Komplex mit NP[\type{str}]$_j$ in \comps. \emph{wurde} selegiert ein Verb oder einen Verbalkomplex mit der Verbform \type{ppp} und den
Argumenten des eingebetteten Verbs. Daher landet NP[\type{str}]$_j$ als erstes Element der
\argstl von \emph{werden}, wo es Nominativ bekommt. Da \emph{werden} finit ist, werden alle \argst-Elemente
auf \comps abgebildet und müssen im Satz realisiert werden. Die Analyse des Verbalkomplexes in
(\ref{bsp-zu-reparieren-versucht-wurde-three}) ist in Abbildung~\vref{fig-zu-reparieren-versucht-wurde} dargestellt.

\begin{figure}
\centering
\begin{forest}
sm edges
[V\feattab{
              \vform \type{fin},\\
              \comps \ibox{1} } 
        [{\ibox{2} V\feattab{
              \vform \type{ppp},\\
              \da    \sliste{ NP[\str]$_i$ },\\
              \comps \ibox{1} }} 
           [{\ibox{3} V\feattab{
              \vform \type{inf},\\
              \subj  \sliste{ NP[\str]$_i$ }, \\ 
              \comps \ibox{1} \sliste{ NP[\str]$_j$ } }} [zu reparieren] ]
           [V\feattab{
              \phon \phonliste{ versucht }\\
              \vform \type{ppp},\\
              \da    \ibox{4},\\
              \comps \ibox{1} $\oplus$ \sliste{ \ibox{3} },\\
              \argst \ibox{1} $\oplus$ \sliste{ \ibox{3} } } 
              [V\feattab{
%              \phon \phonliste{ versuch }\\
              \vform \type{ppp},\\
              \da    \ibox{4} \sliste{ NP[\str]$_i$ },\\
              \argst \ibox{4} $\oplus$ \ibox{1} $\oplus$ \sliste{ \ibox{3} } } [versuch-] ] ] ]
        [V\feattab{
              \vform \type{fin},\\
              \comps \ibox{1} $\oplus$ \sliste{
                \ibox{2} }} [wurde] ]
]
\end{forest}
\caption{\label{fig-zu-reparieren-versucht-wurde}Die Analyse des Fernpassivs als Passivierung eines komplexbildenden Verbs}
\end{figure}

Das Fernpassiv ist auch mit Objektkontrollverben möglich, also Verben, die ein Subjekt,
ein Objekt und eine verbale Projektion nehmen, deren Subjekt mit dem Objekt koreferent ist. Ein Beispiel
ist \emph{erlauben}. (\mex{1}a) zeigt das Verb im Aktiv und ohne Verbalkomplexbildung. Das Objekt von \emph{erlauben} \emph{uns} ist koreferent mit dem
Subjekt von \emph{den Erfolg auszukosten}. (\mex{1}b,c) zeigen, dass das Objekt
von \emph{auszukosten} im Nominativ realisiert werden kann, wenn die Verben einen Verbalkomplex bilden:
\eal
\label{bsp-auskosten-fernpassiv}
\ex 
Sie erlauben uns nicht, den Erfolg auszukosten.
\ex\iw{erlauben}
Keine Zeitung          wird   ihr       zu lesen erlaubt.\footnote{
        Stefan Zweig. \emph{Marie Antoinette}. Leipzig: Insel-Verlag. 1932, S.\,515,
        zitiert nach \citew[\page 309]{Bech55a}.
        Dass dies ein Fall des Fernpassivs ist, wurde von \citet[\page 13]{Askedal88} angemerkt.
}

\ex\iw{auskosten}
Der Erfolg         wurde  uns       nicht auszukosten erlaubt.\footnote{
        \citew[\page 110]{Haider86c}%
}

\label{bsp-auskosten-fernpassiv-haider}
\zl

%\largerpage
\noindent
Das Passiv der Konstruktion ohne Verbalkomplex ist ein unpersönliches Passiv:

\ea
Uns       wurde  erlaubt, den        Erfolg  auszukosten.
\z

\noindent
(\mex{1}) zeigt die \argstv von \stem{erlaub}: \emph{erlauben} nimmt ein Subjekt und ein Dativobjekt. Das Dativobjekt ist mit dem Subjekt des eingebetteten Verbs koindiziert, das heißt, die beiden NPs
haben denselben Index, nämlich $j$.
\ea
\stem{erlaub}:\\
\oneline{%
\avm{
[ arg-st & < !NP[\type{str}]$_i$, NP[\ldat]$_j$!  > \+ \1 \+ < V[\type{inf}, \subj < !NP[\type{str}]$_j$! >, \comps \1 ] > ]
}}
\z
Die Komplemente des eingebetteten Verbs \iboxb{1} werden vom einbettenden Verb übernommen. (\mex{1}) zeigt
die \argstv des jeweiligen höchsten Verbs:
\ea
\label{ex-arg-st-fernpassiv-object-control}
\oneline{%
\begin{tabular}[t]{@{}l@{~}l@{ }l@{}}
a. & \emph{zu lesen erlauben}: & \sliste{ NP[\type{str}]$_i$, NP[\ldat]$_j$, NP[\type{str}]$_k$, V[\comps \sliste{ NP[\type{str}]$_k$ }] }\\[2mm]

b. & \emph{zu lesen erlaubt wird}: & \sliste{ NP[\ldat]$_j$, NP[\type{str}]$_k$, V[\comps \sliste{ NP[\type{str}]$_k$ }] }\\
\end{tabular}
}
\z
(\mex{0}a) zeigt, wie das Objekt von \emph{lesen} angezogen wird, sodass die kombinierte \argst
drei NPs enthält. (\mex{0}b) zeigt die Passivvariante, in der das Subjekt von \emph{erlauben}
unterdrückt ist. Das Ergebnis ist eine \argstl, die mit einer dativischen NP beginnt, einer NP mit lexikalischem Kasus. Da die
erste NP mit strukturellem Kasus Nominativ bekommt und mit dem finiten Verb kongruiert\is{Kongruenz}, macht die Theorie die
richtigen Vorhersagen sogar in so komplexen Situationen wie dem Fernpassiv mit Objektkontrollverben. Die
erste NP mit strukturellem Kasus ist das Subjekt im \ili{Deutsch}en.

%% \subsubsection{Complex passives}

%% \begin{itemize}
%% \item Complex Passives:

%% \ea
%% \gll at Bilen           blev forsøgt repareret\\
%%      that car.\textsc{def} was  tried   repaired\\
%% \glt `that an attempt was made to repair the car'
%% \z

%% \item Raising in passive only.

%% \item \emph{forsøgt} (`to try') does not even take a participle in the active:
%% \eal
%% \ex[]{
%% \gll at   Peter har  forsøgt \emphbf{at} \emphbf{reparere} bilen\\
%%      that Peter has  tried   to repair   car.\textsc{def}\\
%% \glt `that Peter tried to repair the car'
%% }
%% \ex[*]{
%% \gll at   Peter har  forsøgt \emphbf{repareret} bilen\\
%%      that Peter has  tried   repaired car.\textsc{def}\\
%% %\glt `that an attempt was made to repair the car'
%% }
%% \zl

%% %% \item Conclusion: We need special lexical items for passive participles.

%% %% \item analysis of the \ili{German} passive and perfect can be maintained,\\
%% %% compatible with a more general analysis of the passive

%% \end{itemize}

\section{Alternativen}

Wie bei den Abschnitten über Alternativen in den vorigen Kapiteln ist dieser Abschnitt nur für fortgeschrittene
Leser. Es ist nicht notwendig, ihn zu lesen, um den Rest des Buches zu verstehen.

\subsection{Government-\&-Binding-Analysen}
\label{sec-passive-GB}

Die hier übernommene Analyse wurde aus Vorschlägen von Hubert Haider
\citeyearpar{Haider86} entwickelt. Haider hat Analysen im Rahmen von Government \& Binding
\citep{Chomsky81a} entwickelt, und seine Analysen verschiedener Phänomene -- nicht nur des Passivs -- sind weitgehend mit HPSG-Sichtweisen kompatibel und werden
von vielen Forschern übernommen, die im Rahmen von HPSG zum \ili{Deutsch}en arbeiten. Die gängigste Analyse des
Passivs in GB ist jedoch anders. \cites[155--157]{Grewendorf88a}[\page 1311]{Grewendorf93} und
\citet[462--463]{SS88a} haben in ihrem äußerst einflussreichen Lehrbuch \emph{Bausteine syntaktischen Wissens} Chomskys Analyse des Passivs im
\ili{Englisch}en \citep[\page 124]{Chomsky81a} auf das \ili{Deutsch}e übertragen (siehe auch \citealt[\page 180]{Lohnstein2014a} für einen
neueren Vorschlag in diese Richtung in einem Lehrbuch). Diese Analyse basiert auf dem CP/TP/VP-System (siehe Abschnitt~\ref{sec-cp-tp-vp} für eine Diskussion des
Scramblings in diesem System). Die Diskussion dieser Passivanalyse basiert auf \citew[Abschnitt~3.4]{MuellerGT-Eng5}.

%\largerpage[2]
Die Passivanalyse in GB ähnelt der hier vorgeschlagenen Analyse darin, dass sie eine lexikalische Analyse ist:
Der Lexikoneintrag für das Partizip ist insofern besonders, als er dem Akkusativobjekt keinen Kasus zuweist. Es
gibt einen Kasusfilter\is{Kasus!Filter}, der verlangt, dass jede NP in einem Satz Kasus haben muss. Da das Verb keinen
Kasus zuweist, muss die NP, die im Aktiv Akkusativ hätte, ihren Kasus anderswo bekommen. Es gibt zwei
Möglichkeiten, Kasus zu bekommen: Das Subjekt erhält Kasus von (finitem) T\is{Kategorie!funktionale!T}, und
der Kasus der übrigen Argumente kommt von V (\citealp[\page 50]{Chomsky81a}; \citealp[\page
26]{Haider84b}; \citealp[\page 71--73]{FF87a}). Dies wird als Kasusprinzip formuliert:
\begin{principle-break}[Kasusprinzip]\label{Kasusprinzip-GB}
\begin{itemize}
% hier stand `seinem', aber in den Abbildungen gibt es zwei Komplemente
\item V weist seinem Komplement objektiven Kasus (Akkusativ) zu, wenn es strukturellen Kasus trägt.
\item Wenn es finit ist, weist Tense dem Subjekt Kasus zu.
\end{itemize}
\end{principle-break}

\noindent
Abbildung~\ref{Abb-GB-Aktiv} zeigt das Kasusprinzip in Aktion am Beispiel in
(\mex{1}a).\footnote{\label{fn-semantic-role-phrase-boundary}%
Die Abbildung entspricht nicht der \xbar-Theorie in ihrer klassischen Form, da \emph{der Frau}
ein Komplement ist, das mit V$'$ kombiniert wird. In der klassischen \xbar-Theorie müssen alle Komplemente
mit \vnull kombiniert werden. Das führt bei ditransitiven\is{Verb!ditransitives} Strukturen zu einem Problem, da die Strukturen binär sein müssen (siehe \citew{Larson88a} für eine Behandlung von Doppelobjektkonstruktionen).
Außerdem wurde das Verb in den folgenden Abbildungen aus Gründen der Übersichtlichkeit in \vnull belassen. Um
eine wohlgeformte S"=Struktur zu erzeugen, müsste sich das Verb zu seinem Affix in \inull bewegen. Man beachte auch,
dass die Zuweisung der Subjekt-Thetarolle durch das Verb eine Phrasengrenze überschreitet. Dieses Problem
kann gelöst werden, indem man annimmt, dass das Subjekt innerhalb der VP erzeugt wird, dort eine Thetarolle bekommt und
sich dann nach SpecIP bewegt. Ein alternativer Vorschlag war anzunehmen, dass die VP eine semantische Rolle an
SpecIP zuweist \parencites[\page 104--105]{Chomsky81a}[\page 229]{AS83a}.%
}
\eal
\ex 
{}[dass] der Delphin dem Kind den Ball gibt
\ex 
{}[dass] der Ball dem Kind gegeben wird
\zl
%\largerpage[2]
%\fi
\begin{figure}
\hfill
\begin{forest}
sm edges
[TP
  [{NP[nom]}, name=subject [der Delphin, roof]]
  [T\rlap{$'$}
	[VP
		[V\rlap{$'$}
			[{NP[dat]}, name=dobject [dem Kind, roof]]
			[V\rlap{$'$}
				[{NP[acc]},   name=aobject [den Ball, roof]]
				[V ,name=verb    [gib-]]]]]
	[T , name=Infl [-t]]]]
\draw[->,dotted] (Infl.north) .. controls (3.5,-0.1) and (-1.5,0.4)  .. ($(subject.north)+(-.1,.1)$);
\draw[->]        (verb.north) .. controls (2.8,-2.9) and (-.4,.5)   .. ($(subject.north)+(0,.1)$);
\draw[->,dashed] (verb.north) .. controls (2.8,-3.0) and (0.2,-3.05)   .. ($(dobject.north)+(0,.1)$);
%\DrawControl{(2.8,-3)}{1}; \DrawControl{(0.2,-3.05)}{2};
\draw[->,dashed] (verb.north) .. controls (2.3,-4.2) and (1.8,-4.3) .. ($(aobject.north)+(0,.1)$);
%\DrawControl{(2.3,-4.2)}{1}; \DrawControl{(1.8,-4.3)}{2};
%
%\draw (-4,-7) to[grid with coordinates] (4,0.5);
\end{forest}\hfill
\begin{tabular}[b]{ll@{}}
\tikz[baseline]\draw[dotted](0,1ex)--(1,1ex);&nur Kasus\\
\tikz[baseline]\draw(0,1ex)--(1,1ex);&nur Thetarolle\\
\tikz[baseline]\draw[dashed](0,1ex)--(1,1ex);&Kasus und Thetarolle
\\
\\
\end{tabular}
\caption{\label{Abb-GB-Aktiv}Kasus- und Thetarollenzuweisung in Aktivsätzen}
\end{figure}%
%
Die Passivmorphologie von \emph{gegeben} blockiert das Subjekt und absorbiert den strukturellen
Akkusativ. In GB wird angenommen, dass semantische Rollen bestimmten Baumpositionen zugewiesen werden. Diese
Positionen werden in der sogenannten Grundkonfiguration bestimmt, also der Konfiguration, bevor irgendeine
Bewegung und Umorganisation von Bäumen stattfindet. Das Objekt, das im Aktiv Akkusativ bekäme,
erhält im Passiv in seiner Grundposition nur eine semantische Rolle, bekommt aber nicht den
absorbierten Kasus. Daher muss es sich an eine Position bewegen, an der ihm Kasus zugewiesen werden kann
\citep[\page 124]{Chomsky81a}. Abbildung~\ref{Abb-GB-Passiv} zeigt, wie das für Beispiel
(\mex{0}b) funktioniert.
\begin{figure}%[t]
\scalebox{1}{%
\begin{forest}
sm edges
[TP
[{NP[nom]}, name=subject [der Ball$_i$,roof]]
[T\rlap{$'$}
	[VP
		[V\rlap{$'$}
			[{NP[dat]}, name=dobject [dem Kind, roof]]
			[V\rlap{$'$}
				[NP,   name=aobject [\_$_i$]]
				[V ,name=verb [gegeben wir-, roof]]]]]
	[T  ,name=Infl [-\/d]]]]
\draw[->,dotted] (Infl.north) .. controls (2.4,.6)   and (-1.5,-.05) .. ($(subject.north)+(0,.1)$);
%\DrawControl{(2.4,.6)}{1}; \DrawControl{(-1.5,-.05)}{2};
\draw[->,dashed] (verb.north) .. controls (2.2,-3.2)  and (0,-3.0)    .. ($(dobject.north)+(0,.1)$);
%\DrawControl{(2.2,-3.2)}{1}; \DrawControl{(0,-3.0)}{2};
\draw[->]        (verb.north) .. controls (2.0,-4.25) and (1.2,-4.25) .. ($(aobject.north)+(0,.1)$);
%\DrawControl{(2.0,-4.25)}{1}; \DrawControl{(1.2,-4.25)}{2};
%\draw (-3,-7) to[grid with coordinates] (3.6,0.5);
\end{forest}}\hfill
\begin{tabular}[b]{ll@{}}
\tikz[baseline]\draw[dotted](0,1ex)--(1,1ex);&nur Kasus\\
\tikz[baseline]\draw(0,1ex)--(1,1ex);&nur Thetarolle\\
\tikz[baseline]\draw[dashed](0,1ex)--(1,1ex);&Kasus und Thetarolle
\\
\\
\end{tabular}
\caption{\label{Abb-GB-Passiv}Kasus- und Thetarollenzuweisung in Passivsätzen}
\end{figure}%
%\if 0

Diese bewegungsbasierte Analyse funktioniert für das \ili{Englisch}e\il{Englisch} gut, da sich das zugrundeliegende Objekt immer bewegen muss:

\eal
\ex[]{
The dolphin gave [the child] [a ball].
}
\ex[]{
{}[The child] was given [a ball] (by the dolphin).
}
\ex[*]{
It was given [the child] [a ball].
}
\zl
%
(\mex{0}c) zeigt, dass das Füllen der Subjektposition mit einem Expletivum nicht möglich ist, und da
das \ili{Englisch}e ein Subjekt verlangt, muss sich das Objekt also tatsächlich bewegen. \citet[Abschnitt~4.4.3]{Lenerz77} hat jedoch gezeigt, dass eine solche Bewegung im
\ili{Deutsch}en nicht obligatorisch ist. (\mex{1}) veranschaulicht dies:

\eal
\label{ex-passive-German-no-movement}
\ex 
weil der Delphin dem Kind den Ball gab
\ex 
weil    dem        Kind  der        Ball gegeben wurde
\ex 
weil    der        Ball dem        Kind  gegeben wurde
\zl
%\addlines[1]
%\largerpage
Im Vergleich zu (\mex{0}c) ist (\mex{0}b) die unmarkierte Abfolge. \emph{der Ball} in (\mex{0}b) steht
in derselben Position wie \emph{den Ball} in (\mex{0}a), das heißt, es ist keine Bewegung nötig. Nur der Kasus unterscheidet sich.
(\mex{0}c) ist im Vergleich zu (\mex{0}b) jedoch etwas markiert. Wenn man also annähme, dass (\mex{0}c)
die normale Abfolge für Passive ist und (\mex{0}b) daraus durch Bewegung von \emph{dem
  Kind} abgeleitet wird, sollte (\mex{0}b) markierter sein als (\mex{0}c), was den
Fakten widerspricht. Um dieses Problem zu lösen, wurde für Fälle wie (\mex{0}b) eine Analyse mit abstrakter Bewegung vorgeschlagen: Die Elemente bleiben in ihren Positionen, sind aber mit
der Subjektposition verbunden und erhalten ihre Kasusinformation von
dort. \textcites[155--157]{Grewendorf88a}[\page 1311]{Grewendorf93} und \citet[462--463]{SS88a}
nehmen an, dass es ein leeres Expletivpronomen\is{leeres Element}\is{Pronomen!expletives}
% Fanselow81a:152 Infl weist Kasus in die VP zu. Adjazens nicht nötig.
in der Subjektposition von Sätzen wie (\mex{0}b) und in der Subjektposition von Sätzen mit einem
unpersönlichen Passiv\is{Passiv!unpersönliches} wie (\mex{1}) gibt:\footnote{%
	Siehe \citew[\page 11--12]{Koster86a} für eine parallele Analyse des \ili{Niederländisch}en\il{Niederländisch} sowie
	\citew[\page 180]{Lohnstein2014a} für eine bewegungsbasierte Behandlung des Passivs, die für die Analyse des unpersönlichen Passivs ebenfalls ein
        leeres Expletivum involviert.
}
\ea
weil    heute nicht gearbeitet wird
\z
Ein leeres Expletivpronomen ist etwas, das man weder sehen noch hören kann und das keine
Bedeutung trägt. Solche leeren Elemente werden von vielen Forschern abgelehnt, da unklar ist, wie ihre
Existenz von Sprachlernern erworben werden soll. Dafür scheint es notwendig zu sein, reiches angeborenes
sprachliches Wissen anzunehmen, etwas, das heutzutage nicht einmal Chomsky annimmt \citep*{HCF2002a}. Zur Diskussion
dieser Art von leerem Element siehe \citew[Abschnitt~13.1.3 und Kapitel~19]{MuellerGT-Eng5} und \citew[Abschnitt~7]{MuellerHeadless}.

\citet[\page 12]{Koster86a} hat darauf hingewiesen, dass das Passiv im \ili{Englisch}en\il{Englisch} nicht durch die Kasustheorie
abgeleitet werden kann, also durch Kasuslosigkeit und Bewegung in die Spezifikatorposition von TP aufgrund des Kasusfilters, denn wenn man leere Expletivsubjekte sowohl für das \ili{Englisch}e als auch für das \ili{Deutsch}e und \ili{Niederländisch}e\il{Niederländisch} zuließe, dann wäre es möglich,
Analysen wie die folgende in (\mex{1}) zu haben, wobei np ein leeres Expletivum ist:
\ea
np was read the book.
\z
Das Objekt müsste sich nicht in die Subjektposition bewegen und könnte einfach dort bleiben, was den Fakten widerspricht.
Koster nimmt stattdessen an, dass Subjekte im \ili{Englisch}en entweder durch andere Elemente gebunden (also nicht-expletiv) oder lexikalisch gefüllt sind, also
durch sichtbares Material gefüllt.
Daher wäre die Struktur in (\mex{0}) ausgeschlossen, und es wäre sichergestellt, dass \emph{the book} vor dem
finiten Verb platziert werden müsste, sodass die Subjektposition gefüllt ist.
\is{Passiv|)}

Abschließend lässt sich sagen, dass das Passiv nicht durch Bewegung erklärt werden sollte. Chomskys Analyse funktioniert,
aber nur aufgrund der Tatsache, dass das \ili{Englisch}e ein Subjekt verlangt. Es werden zwei Phänomene vermischt, die
getrennt behandelt werden sollten. In der hier vorgeschlagenen Analyse ist das Passiv die Unterdrückung des
Subjekts in der Argumentstrukturliste. Das funktioniert für alle untersuchten Sprachen. Es hängt von der
jeweiligen Sprache ab, ob es ein Subjekt geben muss und ob es in
einer bestimmten Position realisiert werden muss. Im \ili{Englisch}en müssen wir etwas in der Spezifikatorposition haben, im \ili{Deutsch}en sind alle
Argumente auf der \compsl aufgeführt. Es ist keine Bewegung beteiligt. Man beachte, dass das Äquivalent zur
Grundabfolge für die Zuweisung semantischer Rollen die \argstl ist. Die \argstl hat eine bestimmte feste Abfolge. Diese Tatsache
kann für das Linking genutzt werden, da klar ist, an welcher Position der Liste welches Argument
repräsentiert ist. Passivsätze enthalten ein passiviertes Verb. Das passivierte Verb ist mit einem Verbstamm
mit einer \argstl verbunden, die der Aktivform entspricht. Diese \argstl ist also irgendwie für die
Zuweisung semantischer Rollen zugänglich, aber die Analyse des Passivsatzes umfasst nicht die Analyse eines
Aktivsatzes und irgendeine Bewegung. Nach allem, was bisher über die menschliche Sprachverarbeitung bekannt ist, ist dies
aus psycholinguistischer Sicht der richtige Ansatz \citep[Abschnitt~3.2]{Wasow2021a}. 

\section{Zusammenfassung}

Zusammenfassend lässt sich sagen, dass dieses Kapitel eine einheitliche Analyse des Passivs im \ili{Dänisch}en,
\ili{Englisch}en, \ili{Deutsch}en und \ili{Isländisch}en liefert. Die Lexikonregel erfasst sowohl die morphologischen als auch die
analytischen Passive. Das erste Element auf der \argstl wird unterdrückt, und eine relationale Beschränkung
\texttt{promote} befördert jede NP mit strukturellem Kasus. Die Sprachen unterscheiden sich darin, welche Kasus sie verwenden und
welche Kasus strukturell/lexikalisch sind. Das \ili{Dänisch}e fügt Expletiva ein, um unpersönliche Passive zu erlauben und
das Bedürfnis nach einem Subjekt zu erfüllen. Diese Expletiveinfügung erfolgt in der \argst-Abbildung, wenn
Argumente von \argst auf die Valenzlisten abgebildet werden.

Die SVO-Sprachen scheinen unterschiedliche Einträge für die Perfekt-/Passivpartizipien zu erfordern, aber \citegen{Haider86} Passivanalyse für das \ili{Deutsch}e, die nur eine Partizipform für sowohl Perfekt als auch Passiv verwendet, kann beibehalten werden.

\questions{
\begin{enumerate}
\item Welche Tests für Subjekteigenschaft kennen Sie?
\item Funktionieren diese Tests für alle \ili{germanisch}en Sprachen?
\item Inwiefern unterscheidet sich das \ili{Deutsch}e vom \ili{Isländisch}en in Bezug auf Subjekte?
\item Was ist struktureller Kasus? Was ist lexikalischer Kasus?

\item Was ist ein unpersönliches Passiv?
\item Hat das \ili{Isländisch}e unpersönliche Passive?

\end{enumerate}

}

%\largerpage
\exercises{
\begin{enumerate}
\item Welche NPs in (\mex{1}) haben strukturellen und welche lexikalischen Kasus?
\eal
\ex 
Der        Junge lacht.
\ex 
Mich friert.
\ex 
Er zerstört die Umwelt.
\ex 
Das dauert ein ganzes Jahr.
\ex 
Er hat nur einen Tag dafür gebraucht.
\ex 
Er denkt an den morgigen Tag.
\zl

\item Geben Sie \argst-Listen für die folgenden Verben an:
\eal
\ex show, eat, meet \english
\ex zeigen, essen, begegnen, treffen \german
%\ex vise `show', spise `eat', møde `meet' \danish
%\ex sýna `show', eta `eat', mæta `meet', hitta `meet' \icelandic
\zl
Wenn Sie sich beim Kasus unsicher sind, können Sie das
  Wiktionary verwenden: \url{https://de.wiktionary.org/}.

\item Zeichnen Sie den Analysebaum für den folgenden Satz:
\ea
that the box was opened
\z
Geben Sie die Valenzmerkmale (\spr und \comps) und die Wortartinformation an. Sie können
die NP mit einem Dreieck abkürzen.

%% \item Draw the analysis tree for the following sentence:
%% \ea
%% \gll dass der Kasten geöffnet wurde\\
%%      that the box    opened was\\
%% \glt `that the box opened was'
%% \z
%% Please provide valence features (\comps only) and part of speech information. You may abbreviate
%% the NP using a triangle.

\end{enumerate}

}

%      <!-- Local IspellDict: de_DE -->
